\documentclass[11pt,oneside]{article}

\usepackage[T1]{fontenc}
\usepackage{lmodern}
\usepackage{microtype}
\usepackage[margin=1in,headheight=14pt]{geometry}
\usepackage{amsmath,amssymb,amsthm,mathtools}
\usepackage{booktabs,longtable,tabularx,array}
\usepackage{enumitem}
\usepackage{xcolor}
\usepackage{xurl}
\usepackage{hyperref}
\usepackage{fancyhdr}
\usepackage{listings}
\usepackage{needspace}

\definecolor{provedgreen}{RGB}{17,102,72}
\definecolor{correctionred}{RGB}{150,39,39}
\definecolor{classicalblue}{RGB}{35,82,130}
\definecolor{openviolet}{RGB}{102,55,145}
\definecolor{codegray}{RGB}{247,247,247}
\definecolor{rulegray}{RGB}{95,102,110}

\hypersetup{
  colorlinks=true,
  linkcolor=classicalblue,
  citecolor=provedgreen,
  urlcolor=classicalblue,
  pdftitle={Alaoglu--Erdos Progress Report 7: Rotational Regularization and the Exact Quadratic-Type Barrier},
  pdfauthor={Research progress report},
  pdfsubject={Alaoglu--Erdos conjecture, entire interpolation, Carlson--Carleman uniqueness, difference algebra}
}
\urlstyle{same}

\pagestyle{fancy}
\fancyhf{}
\lhead{Alaoglu--Erd\H{o}s progress report 7}
\rhead{Quadratic type and translation rank}
\cfoot{\thepage}
\renewcommand{\headrulewidth}{0.4pt}

\setlength{\parindent}{0pt}
\setlength{\parskip}{0.55em}
\setlength{\emergencystretch}{4em}
\setlist[itemize]{leftmargin=2em,itemsep=0.28em,topsep=0.4em}
\setlist[enumerate]{leftmargin=2.2em,itemsep=0.28em,topsep=0.4em}

\newtheorem{theorem}{Theorem}[section]
\newtheorem{proposition}[theorem]{Proposition}
\newtheorem{lemma}[theorem]{Lemma}
\newtheorem{corollary}[theorem]{Corollary}
\newtheorem{conjecture}[theorem]{Conjecture}
\newtheorem{question}[theorem]{Question}
\newtheorem{target}[theorem]{Research target}
\theoremstyle{definition}
\newtheorem{definition}[theorem]{Definition}
\newtheorem{example}[theorem]{Example}
\newtheorem{warning}[theorem]{Warning}
\newtheorem{experiment}[theorem]{Computational experiment}
\theoremstyle{remark}
\newtheorem{remark}[theorem]{Remark}

\newcommand{\N}{\mathbb N}
\newcommand{\Nzero}{\mathbb N_0}
\newcommand{\Z}{\mathbb Z}
\newcommand{\Q}{\mathbb Q}
\newcommand{\R}{\mathbb R}
\newcommand{\C}{\mathbb C}
\newcommand{\Abar}{\overline{\mathbb Q}}
\newcommand{\Sbeta}{\mathcal S_\beta}
\newcommand{\Zset}{\mathcal Z}
\newcommand{\Eclass}{\mathcal E}
\newcommand{\Tclass}{\mathcal T}
\newcommand{\Vclass}{\mathcal V}
\newcommand{\M}{\mathcal M}
\newcommand{\ee}{\mathrm e}
\newcommand{\ord}{\operatorname{ord}}
\newcommand{\trdeg}{\operatorname{trdeg}}
\newcommand{\Span}{\operatorname{span}}
\newcommand{\supp}{\operatorname{supp}}
\newcommand{\dist}{\operatorname{dist}}
\newcommand{\type}{\operatorname{type}}
\newcommand{\DC}{\mathfrak D_{\!C}}
\newcommand{\Cbeta}{C_\beta}
\newcommand{\TP}{T_{\mathrm P}(\beta)}
\newcommand{\Tsh}[1]{T_{#1}}
\newcommand{\Del}[1]{\Delta_{#1}}
\newcommand{\statusproved}{\textcolor{provedgreen}{\textbf{[proved here]}}}
\newcommand{\statuscorrection}{\textcolor{correctionred}{\textbf{[correction]}}}
\newcommand{\statusclassical}{\textcolor{classicalblue}{\textbf{[classical input]}}}
\newcommand{\statusopen}{\textcolor{openviolet}{\textbf{[open target]}}}
\newcommand{\statuscomputed}{\textcolor{openviolet}{\textbf{[exact computation]}}}

\newenvironment{statusbox}[1]
{\begin{center}\begin{minipage}{0.94\textwidth}\raggedright\hrule\vspace{0.45em}\textbf{Status:} #1\par\vspace{0.35em}}
{\vspace{0.35em}\hrule\end{minipage}\end{center}}

\lstdefinestyle{pythonstyle}{
  language=Python,
  basicstyle=\ttfamily\scriptsize,
  backgroundcolor=\color{codegray},
  frame=single,
  rulecolor=\color{rulegray},
  numbers=left,
  numberstyle=\tiny\color{rulegray},
  numbersep=7pt,
  showstringspaces=false,
  breaklines=true,
  columns=fullflexible,
  keepspaces=true,
  tabsize=4
}

\begin{document}
\pagenumbering{gobble}

\begin{titlepage}
\centering
\vspace*{0.65cm}
{\Huge\bfseries The Alaoglu--Erd\H{o}s Conjecture\par}
\vspace{0.22cm}
{\LARGE\bfseries Progress Report 7\par}
\vspace{0.75cm}
{\LARGE Rotational Regularization, the Exact Quadratic-Type Barrier,\par}
\vspace{0.12cm}
{\LARGE and a New Translation-Rank Route\par}
\vspace{0.9cm}
{\large A corrective and forward-looking continuation of\par}
\vspace{0.1cm}
{\large \texttt{alaoglu-erdos-unified-report-6.zip}\par}
\vfill
\begin{minipage}{0.92\textwidth}
\small
\textbf{Bottom line.}
The conjecture
\[
  2^x,3^x\in\Z\quad\Longrightarrow\quad x\in\Z
\]
is not proved here.  A substantial analytic part of the previous frontier is, however,
replaced by an exact theorem.  Conditional on a counterexample with least generator
$\beta$, the solution monoid is $\Sbeta=\Nzero+\beta\Nzero$ and has quadratic density
$1/(2\beta)$.  The preceding report asserted that every entire function vanishing on
$\Sbeta$ has infinite quadratic type.  That assertion is false.  The explicit product
\[
 H_\beta(z)=z\prod_{s\in\Sbeta\setminus\{0\}}
                 \left(1-\frac{z^4}{s^4}\right)
\]
is entire, vanishes on $\Sbeta$, and has the exact finite quadratic type
\[
  \tau_2(H_\beta)=\frac{\pi}{2\beta}.
\]
A Carlson--Carleman half-plane argument proves that no nonzero entire function vanishing
on $\Sbeta$ can have smaller type.  In fact the complete finite vanishing spectrum is
$[\pi/(2\beta),\infty)$.

This exact threshold sharpens the Pila interpolation programme.  It proves that
subcritical interpolation on $\Sbeta$ is unique, quantifies how many values or templates a
Pila-small integer-valued function must use, excludes polynomial coordinate decoders and
finite constant-coefficient recurrences, and shows that no null-function perturbation can
come remotely close to Pila's sufficient threshold.  A second new route comes from Pila's
multi-function theorem: a sufficiently small integer-valued interpolant has a quantitatively
finite algebraic translation rank under the two commuting shifts $1$ and $\beta$.  Thus a
proof may be sought through the classification of finite-order, transformally algebraic
entire functions for two incommensurable shifts.  The linear and polynomial-decoder cases
are closed here; the genuinely nonlinear difference-algebraic case is the new surviving
target.
\end{minipage}
\vfill
{\large August 24, 2026\par}
\end{titlepage}

\pagenumbering{roman}

\begin{abstract}
Assume, for contradiction, that the Alaoglu--Erd\H{o}s conjecture fails.  The exact
conditional structure developed in the predecessor report supplies a unique least
nonintegral solution $\beta$ and the free additive monoid
$\Sbeta=\{n+k\beta:n,k\in\Nzero\}$.  Its counting function is
$R^2/(2\beta)+O_\beta(R)$.  This report solves the associated zero-set problem at order
$2$.  For every integer $m\ge3$ the rotationally regularized product
\[
 H_{m,\beta}(z)=z\prod_{s\in\Sbeta\setminus\{0\}}
                    \left(1-\frac{z^m}{s^m}\right)
\]
is entire and satisfies
\[
 \log M(H_{m,\beta},r)
 =\frac{\pi}{2\beta\sin(2\pi/m)}r^2+O_{\beta,m}(r+\log r).
\]
The unique optimal integer rotation order is $m=4$.  Conversely, if a nonzero entire
function of finite quadratic type vanishes on $\Sbeta$, composing with the principal
square root turns it into a right-half-plane function of exponential type with zeros of
linear density $1/(2\beta)$.  Carleman's formula gives the sharp lower bound
$\tau_2\ge\pi/(2\beta)$.  Hence the exact finite vanishing spectrum is
$[\pi/(2\beta),\infty)$, and the report's earlier infinite-type proposition is corrected.

Let $\Tclass_\beta$ be the set of finite quadratic types of entire functions integer-valued
on $\Sbeta$.  We prove
\[
 \{0\}\cup\left[\frac{\pi}{2\beta},\infty\right)\subseteq\Tclass_\beta,
\]
whereas Pila's theorem and the algebraic-dependence growth lemma exclude the interval
\[
  \left(0,\frac{1}{832\beta^2\log2\log3}\right].
\]
The only unresolved type range is therefore a finite
middle interval, not the existence of positive finite type.  A Carleman-density covering
theorem yields quantitative template, alphabet, and coincidence bounds.  It implies that
a function at Pila's threshold must use at least approximately
$416\pi\beta\log2\log3$ distinct states in any finite-value decoder unless the decoder
collapses to a global identity.  We also prove that a subcritical polynomial decoder
$g(n+k\beta)=P(n,k)$ must factor through $n+k\beta$; for algebraic-coefficient $P$ it is
constant.  Two constant-coefficient recurrences in the coordinate directions force an
exponential polynomial and hence quadratic type zero.

Finally, specializing Pila's multi-function theorem gives an explicit weighted-simplex
criterion for several integer-valued entire functions.  Applied to translates
$g(z+h)$, $h\in\Sbeta$, it bounds the relative transcendence degree of the full translation
orbit whenever $2\beta\tau_2(g)<29/75$.  This identifies a new proof programme in partial
difference algebra.  Every assertion is labelled as proved here, classical input, exact
computation, correction, or open target; no competitive announcement is used as a
mathematical premise.
\end{abstract}

\tableofcontents
\newpage
\pagenumbering{arabic}

\section{Scope, inherited structure, and the new frontier}
\label{sec:scope}

\subsection{The conjecture and the conditional monoid}

Write
\[
 \vartheta=\frac{\log3}{\log2}.
\]
The Alaoglu--Erd\H{o}s conjecture is equivalently the assertion that
$m^\vartheta\in\Z$ for $m\in\N$ forces $m$ to be a power of two.  It is also the special
positive-integral $2\times2$ case of the four exponentials conjecture obtained from the two
$\Q$-independent pairs $(1,x)$ and $(\log2,\log3)$.

The present report does not repeat the large structural corpus of the predecessor~\cite{Unified6}.  It uses
the following conditional package as input.

\begin{itemize}
\item If a nonintegral solution exists, there is a unique least one, denoted $\beta>0$.
\item The complete solution set is the free additive monoid
\[
 \Sbeta=\{n+k\beta:n,k\in\Nzero\}.
\]
\item The representation $n+k\beta$ is unique because $\beta$ is irrational; indeed
$\beta$ is transcendental by Gelfond--Schneider.
\item Every $s\in\Sbeta$ satisfies $2^s,3^s\in\Z$.
\item The predecessor's kernel-checked finite region gives $\beta\ge14$; its wider
software scan gives $\beta>27$ under the stated computational trust boundary.
\end{itemize}

Only the first four items are needed for the proofs below.  The numerical lower bounds on
$\beta$ are used solely to illustrate the size of the resulting constants.

\subsection{What is corrected}

The predecessor report contains the following claim:

\begin{quote}
Every entire function vanishing on all of $\Sbeta$ has infinite quadratic type.
\end{quote}

Its supporting calculation correctly detects that the genus-two canonical product having
\emph{exactly} the positive zeros $\Sbeta$ has growth of order $r^2\log r$ in a bad
direction.  The logical leap from that specific product to all entire functions vanishing
on $\Sbeta$ is incorrect.  At integer order, angular Lindel\"of cancellation matters.
Additional zeros can cancel the divergent reciprocal-square moment and reduce the maximum
modulus from $r^2\log r$ to exact finite quadratic type.  The fourfold rotational product
$H_\beta$ on the title page does precisely this.

\begin{statusbox}{\statuscorrection.  The counterexample is explicit, its convergence is
absolute on compact sets, and its exact type is proved in Section~\ref{sec:rotprod}.}
\end{statusbox}

\subsection{Main theorems of this report}

For an entire function $f$, set
\[
 M(f,r)=\max_{|z|=r}|f(z)|,
 \qquad
 \tau_2(f)=\limsup_{r\to\infty}\frac{\log M(f,r)}{r^2}
 \in[0,\infty].
\]
The report establishes the following new statements.

\begin{enumerate}
\item \textbf{Exact rotational products.}  For every integer $m\ge3$,
\[
 \tau_2(H_{m,\beta})
 =\frac{\pi}{2\beta\sin(2\pi/m)}.
\]
The minimum over $m$ is attained uniquely at $m=4$.

\item \textbf{Sharp zero-set threshold.}  If $f\not\equiv0$ is entire, vanishes on
$\Sbeta$, and has finite quadratic type, then
\[
 \tau_2(f)\ge \Cbeta:=\frac{\pi}{2\beta}.
\]
Equality is attained by $H_\beta=H_{4,\beta}$.

\item \textbf{Complete vanishing spectrum.}
\[
 \Vclass_\beta
 :=\{\tau_2(f):f\not\equiv0,\ f|_{\Sbeta}=0,\ \tau_2(f)<\infty\}
 = [\Cbeta,\infty).
\]

\item \textbf{Corrected integer-valued spectrum.}  If
\[
 \Tclass_\beta
 :=\{\tau_2(g):g\text{ entire},\ g(\Sbeta)\subseteq\Z,
                       \ \tau_2(g)<\infty\},
\]
then
\[
 \{0\}\cup[\Cbeta,\infty)\subseteq\Tclass_\beta.
\]
Pila's theorem conditionally excludes a small interval immediately above zero.  Thus the
unknown range is finite.

\item \textbf{Carleman-density covering.}  If finitely many zero sets of entire functions
cover $\Sbeta$, the sum of their quadratic types is at least $\Cbeta$.  This yields exact
finite-template and finite-alphabet complexity lower bounds.

\item \textbf{Decoder rigidity.}  Below $\Cbeta$, polynomial coordinate decoders and two
constant-coefficient recurrences collapse to type zero.  At Pila's much smaller threshold,
any finite-state description must have a state count on the order of $10^4$ already at
$\beta=14$.

\item \textbf{Multi-function and translation-rank gate.}  Pila's general theorem for
several additional functions gives a weighted-simplex criterion.  Applied to the translates
of one interpolant, it forces finite transformal transcendence degree under the shifts
$1$ and $\beta$.  This is the principal new open route.
\end{enumerate}

\subsection{What is not claimed}

No theorem below excludes the middle type interval
\[
 \left(\frac{1}{832\beta^2\log2\log3},\frac{\pi}{2\beta}\right).
\]
No construction in that interval is supplied.  The conjecture therefore remains open.  The
progress is instead an exact correction, a solved extremal problem, several new rigidity
theorems, and a narrower difference-algebraic target.

\section{Counting the conditional solution monoid}
\label{sec:counting}

Throughout the rest of the report we work conditionally: assume a counterexample exists and
let $\beta$ be its canonical least generator.  Put
\[
 \Sbeta=\{n+k\beta:n,k\in\Nzero\},
 \qquad
 N_\beta(R)=\#\bigl(\Sbeta\cap(0,R]\bigr).
\]

\begin{proposition}[Quadratic counting law]
\label{prop:counting}
As $R\to\infty$,
\[
 N_\beta(R)=\frac{R^2}{2\beta}+O_\beta(R).
\]
More explicitly, if $K=\lfloor R/\beta\rfloor$, then
\[
 N_\beta(R)=
 \sum_{k=0}^{K}\bigl(\lfloor R-k\beta\rfloor+1\bigr)-1.
\]
\end{proposition}

\begin{proof}
Irrationality of $\beta$ gives unique coordinates.  For a fixed $k$, the admissible values
of $n$ are $0\le n\le R-k\beta$, hence the displayed exact sum.  Replacing every floor by
its argument creates an error bounded by $K+1=O_\beta(R)$, and
\[
 \sum_{k=0}^{K}(R-k\beta+1)
 =(K+1)(R+1)-\frac{\beta K(K+1)}2
 =\frac{R^2}{2\beta}+O_\beta(R).
\]
Removing the origin changes the result by one.
\end{proof}

Set
\[
 d_\beta:=\frac1{2\beta}.
\]
Thus $N_\beta(R)=d_\beta R^2+O_\beta(R)$.

\begin{corollary}[Reciprocal-power sums]
\label{cor:reciprocal}
As $R\to\infty$,
\[
 \sum_{\substack{s\in\Sbeta\\0<s\le R}}\frac1{s^2}
 =\frac1\beta\log R+O_\beta(1)
 =2d_\beta\log R+O_\beta(1).
\]
For every $m>2$,
\[
 \sum_{s\in\Sbeta\setminus\{0\}}\frac1{s^m}<\infty.
\]
\end{corollary}

\begin{proof}
Apply Stieltjes partial summation to $N_\beta(t)$.  For the critical exponent,
\[
 \sum_{s\le R}s^{-2}
 =\frac{N_\beta(R)}{R^2}+2\int_{1^-}^{R}\frac{N_\beta(t)}{t^3}\,dt
 =2d_\beta\log R+O_\beta(1).
\]
For $m>2$, the corresponding integral has integrand $O(t^{1-m})$ and converges.
\end{proof}

\begin{corollary}[The squared sequence]
\label{cor:squared}
Let
\[
 \Lambda_\beta=\{s^2:s\in\Sbeta\setminus\{0\}\}.
\]
Its counting function satisfies
\[
 N_{\Lambda_\beta}(T)=d_\beta T+O_\beta(\sqrt T).
\]
\end{corollary}

\begin{proof}
This is Proposition~\ref{prop:counting} with $R=\sqrt T$.
\end{proof}

The passage from quadratic density on a ray to linear density after squaring is the source
of the exact constant $\pi d_\beta$ below.

\section{Why the exact-zero canonical product has infinite type}
\label{sec:audit}

The erroneous proposition in the predecessor report arose from a true calculation applied
to the wrong class of functions.  It is useful to preserve the true part.

Let
\[
 E_2(w)=(1-w)\exp\left(w+\frac{w^2}{2}\right)
\]
be the genus-two primary factor, and define the canonical product with \emph{exactly} the
prescribed positive zeros
\[
 P_\beta(z)=z\prod_{s\in\Sbeta\setminus\{0\}}E_2(z/s).
\]
The product converges because $\sum s^{-3}<\infty$.

\begin{proposition}[The one-ray exact-zero product is of infinite quadratic type]
\label{prop:oneray}
Along the negative real axis,
\[
 \log|P_\beta(-r)|\ge d_\beta r^2\log r-O_\beta(r^2).
\]
Consequently $\tau_2(P_\beta)=+\infty$.
\end{proposition}

\begin{proof}
For $x\ge0$,
\[
 \log E_2(-x)=\log(1+x)-x+\frac{x^2}{2}.
\]
The expression is nonnegative because its derivative is $x^2/(1+x)$.  Hence the factors
with $s>r$ contribute nonnegatively.  For $s\le r$, discard the nonnegative logarithm and
sum the remaining terms:
\[
 \begin{aligned}
 \log|P_\beta(-r)|
 &\ge \log r-r\sum_{s\le r}\frac1s
       +\frac{r^2}{2}\sum_{s\le r}\frac1{s^2}.
 \end{aligned}
\]
Partial summation gives $\sum_{s\le r}s^{-1}=2d_\beta r+O_\beta(\log r)$, while
Corollary~\ref{cor:reciprocal} gives
$\sum_{s\le r}s^{-2}=2d_\beta\log r+O_\beta(1)$.  Substitution proves the claim.
\end{proof}

\begin{warning}[The precise logical failure]
Proposition~\ref{prop:oneray} concerns the canonical product whose zero divisor is exactly
$\Sbeta$.  An arbitrary entire function that vanishes on $\Sbeta$ may have additional
zeros.  At integer order two, those extra zeros can cancel the divergent reciprocal-square
angular moment.  It is therefore false that the exact-zero product minimizes maximum-modulus
growth over all zero divisors containing $\Sbeta$.
\end{warning}

The next section gives the elementary regularization.  Its factor has no quadratic term at
all:
\[
 \log\left(1-\frac{z^4}{s^4}\right)
 =-\frac{z^4}{s^4}+O\left(\frac{|z|^8}{s^8}\right).
\]
Thus the divergent sum $\sum s^{-2}$ never appears.

\section{Rotational regularization and exact type}
\label{sec:rotprod}

\subsection{The cyclotomic products}

For an integer $m\ge3$, define
\begin{equation}
 H_{m,\beta}(z)
 :=z\prod_{s\in\Sbeta\setminus\{0\}}
          \left(1-\frac{z^m}{s^m}\right).
 \label{eq:Hm}
\end{equation}
Each factor adds the full $m$-fold root-of-unity orbit of $s$ to the zero set.  In
particular, every positive $s\in\Sbeta$ remains a zero.

\begin{lemma}[Local uniform convergence]
\label{lem:convergence}
For every integer $m\ge3$, the product in \eqref{eq:Hm} converges locally uniformly and
defines an entire function.
\end{lemma}

\begin{proof}
Corollary~\ref{cor:reciprocal} gives $\sum_s s^{-m}<\infty$.  On a compact disk
$|z|\le R$, the series $\sum_s|z|^m/s^m$ converges uniformly.  The standard infinite-product
criterion applies.
\end{proof}

\begin{lemma}[Exact maximum modulus]
\label{lem:maxmod}
For every $r>0$,
\[
 M(H_{m,\beta},r)
 =r\prod_{s\in\Sbeta\setminus\{0\}}
       \left(1+\frac{r^m}{s^m}\right).
\]
Equality is attained at every ray satisfying $z^m=-r^m$.
\end{lemma}

\begin{proof}
For $|z|=r$,
\[
 \left|1-\frac{z^m}{s^m}\right|
 \le1+\frac{r^m}{s^m}.
\]
All these inequalities become equalities simultaneously when $z^m=-r^m$.
\end{proof}

\begin{theorem}[Exact type of the rotational products]
\label{thm:Hmtype}
For every integer $m\ge3$,
\begin{equation}
 \log M(H_{m,\beta},r)
 =\pi d_\beta\csc\left(\frac{2\pi}{m}\right)r^2
   +O_{\beta,m}(r+\log r).
 \label{eq:Hmasymp}
\end{equation}
Equivalently,
\begin{equation}
 \tau_2(H_{m,\beta})
 =\frac{\pi}{2\beta\sin(2\pi/m)}.
 \label{eq:Hmtype}
\end{equation}
\end{theorem}

\begin{proof}
By Lemma~\ref{lem:maxmod} and Stieltjes integration,
\[
 \begin{aligned}
 \log M(H_{m,\beta},r)-\log r
 &=\int_{0^+}^{\infty}\log\left(1+\frac{r^m}{t^m}\right)\,dN_\beta(t)\\
 &=m r^m\int_{0^+}^{\infty}
       \frac{N_\beta(t)}{t(t^m+r^m)}\,dt.
 \end{aligned}
\]
Insert $N_\beta(t)=d_\beta t^2+O_\beta(t+1)$.  The main term is
\[
 m d_\beta r^m\int_0^\infty\frac{t}{t^m+r^m}\,dt
 =m d_\beta r^2\int_0^\infty\frac{u}{1+u^m}\,du.
\]
The beta-integral identity
\[
 \int_0^\infty\frac{u^{a-1}}{1+u^m}\,du
 =\frac{\pi}{m}\csc\left(\frac{\pi a}{m}\right)
 \qquad(0<a<m)
\]
with $a=2$ gives the coefficient in \eqref{eq:Hmasymp}.  The $O(t)$ part of the counting
error contributes $O_{\beta,m}(r)$; the bounded part contributes $O_{\beta,m}(\log r)$.
\end{proof}

\begin{corollary}[The optimal rotational order]
\label{cor:m4}
Among integers $m\ge3$, the type in \eqref{eq:Hmtype} is minimized uniquely at $m=4$.
Thus
\begin{equation}
 H_\beta(z):=H_{4,\beta}(z)
 =z\prod_{s\in\Sbeta\setminus\{0\}}
       \left(1-\frac{z^4}{s^4}\right)
 \label{eq:H4}
\end{equation}
satisfies
\begin{equation}
 \tau_2(H_\beta)=\Cbeta:=\frac{\pi}{2\beta}.
 \label{eq:Cbeta}
\end{equation}
\end{corollary}

\begin{proof}
For $m\ge3$, the quantity $\sin(2\pi/m)$ is at most $1$, with equality exactly for
$m=4$.
\end{proof}

\begin{remark}[Why four rotations are natural]
The exponent of convergence of $\Sbeta$ is two.  A genus-zero cyclotomic factor must begin
at degree strictly larger than two, so $m\ge3$.  The type constant is controlled by the
angle of a maximum-modulus ray.  The fourth roots place that ray at $\pi/4$, exactly the
balanced direction for the squaring map used in the sharp lower bound.
\end{remark}

\begin{statusbox}{The convergence, exact maximum-modulus identity, asymptotic, and optimality
are \statusproved.  They provide an explicit counterexample to the predecessor's
infinite-type proposition.}
\end{statusbox}

\section{A sharp quadratic Carlson--Carleman theorem}
\label{sec:carleman}

The product $H_\beta$ gives an upper bound for the least possible type of a function
vanishing on $\Sbeta$.  The lower bound comes from a half-plane zero-density theorem.  The
proof is a quadratic-order version of Carlson's theorem: square the variable, then apply
Carleman's formula in a half-plane.

\subsection{Carleman's half-plane inequality}

We record only the form needed here.  A derivation from Green's identity appears in
Appendix~\ref{app:carleman}.

\begin{lemma}[Carleman inequality]
\label{lem:carleman}
Let $F$ be holomorphic in the open right half-plane and suppose that for every
$\varepsilon>0$,
\[
 \log^+|F(w)|\le(\sigma+\varepsilon)|w|+O_\varepsilon(1)
 \qquad(\Re w>0).
\]
Let $\{\lambda_j\}$ be the positive real zeros of $F$, counted with multiplicity.  Then,
for boundary-regular radii $R\to\infty$,
\begin{equation}
 \sum_{1<\lambda_j<R}
 \left(\frac1{\lambda_j}-\frac{\lambda_j}{R^2}\right)
 \le \frac{\sigma+\varepsilon}{\pi}\log R+O_{F,\varepsilon}(1).
 \label{eq:carlemanineq}
\end{equation}
The same asymptotic inequality holds for arbitrary $R$ after harmless one-sided limiting.
\end{lemma}

\begin{remark}
For zeros at all positive integers, the left side is $\log R+O(1)$, so
Lemma~\ref{lem:carleman} gives $\sigma\ge\pi$.  This is the classical Carlson constant.
The modern uniqueness-set literature formulates the same phenomenon through angular density
and trigonometrically convex indicators; see Grishin--Sodin and the recent series of
Kononova cited in Section~\ref{sec:literature}.
\end{remark}

\subsection{The general ray-density theorem}

\begin{theorem}[Quadratic ray-density lower bound]
\label{thm:raydensity}
Let $A\subset(0,\infty)$ be discrete and suppose
\[
 N_A(r):=\#(A\cap(0,r])=d r^2+o(r^2)
\]
for some $d>0$.  If $f\not\equiv0$ is entire, $f|_A=0$, and $\tau_2(f)<\infty$, then
\begin{equation}
 \tau_2(f)\ge\pi d.
 \label{eq:raylower}
\end{equation}
If the counting error is $O(r)$, the proof gives an $O(1)$ error after the leading
$d\log R$ Carleman term.
\end{theorem}

\begin{proof}
Choose the principal square root on the right half-plane and put
\[
 F(w)=f(\sqrt w).
\]
If necessary, divide $f$ by a power of $z$ so that behavior at the boundary point $w=0$
does not enter.  The function $F$ is holomorphic for $\Re w>0$.  By the definition of
quadratic type, for every $\varepsilon>0$,
\[
 \log^+|F(w)|
 \le(\tau_2(f)+\varepsilon)|w|+O_\varepsilon(1).
\]
Its positive real zeros include
\[
 \Lambda_A=\{a^2:a\in A\},
\]
whose counting function is $N_{\Lambda_A}(T)=dT+o(T)$.

Apply Lemma~\ref{lem:carleman}.  Partial summation gives
\[
 \begin{aligned}
 \sum_{\substack{a\in A\\1<a^2<R}}
 \left(\frac1{a^2}-\frac{a^2}{R^2}\right)
 &=d\log R+o(\log R).
 \end{aligned}
\]
Indeed,
\[
 \sum_{\lambda<R}\frac1\lambda
 =\frac{N_{\Lambda_A}(R)}R+
   \int_{1^-}^{R}\frac{N_{\Lambda_A}(t)}{t^2}\,dt
 =d\log R+O(1)+o(\log R),
\]
whereas
\[
 \frac1{R^2}\sum_{\lambda<R}\lambda
 =\frac1{R^2}\left(RN_{\Lambda_A}(R)-
   \int_{1^-}^{R}N_{\Lambda_A}(t)\,dt\right)
 =\frac d2+o(1).
\]
Comparing with \eqref{eq:carlemanineq}, dividing by $\log R$, and letting
$\varepsilon\downarrow0$ yields $d\le\tau_2(f)/\pi$.
\end{proof}

\subsection{The exact threshold for the solution monoid}

\begin{theorem}[Sharp quadratic zero-set threshold]
\label{thm:sharpthreshold}
If $f\not\equiv0$ is entire, $f(s)=0$ for all $s\in\Sbeta$, and
$\tau_2(f)<\infty$, then
\begin{equation}
 \tau_2(f)\ge\Cbeta=\frac{\pi}{2\beta}.
 \label{eq:sharpthreshold}
\end{equation}
The constant is sharp, attained by $H_\beta$.
\end{theorem}

\begin{proof}
Proposition~\ref{prop:counting} gives density $d_\beta=1/(2\beta)$, so
Theorem~\ref{thm:raydensity} gives the lower bound.  Corollary~\ref{cor:m4} gives equality.
\end{proof}

\begin{corollary}[Subcritical uniqueness]
\label{cor:uniqueness}
If entire functions $f$ and $g$ satisfy
\[
 \max\{\tau_2(f),\tau_2(g)\}<\Cbeta
\]
and agree on $\Sbeta$, then $f\equiv g$.
\end{corollary}

\begin{proof}
The difference $f-g$ has type at most the maximum of the two types and vanishes on
$\Sbeta$.  Theorem~\ref{thm:sharpthreshold} forces it to be zero.
\end{proof}

\begin{remark}[A true minimal-growth statement]
The false assertion ``every vanishing function has infinite type'' is replaced by a sharp
and complete statement: $\Sbeta$ is a uniqueness set for the open quadratic-type ball
$\tau_2<\pi/(2\beta)$, and it ceases to be a uniqueness set exactly at the boundary.
\end{remark}

\begin{statusbox}{Theorem~\ref{thm:raydensity} is \statusproved from the classical
Carleman formula \statusclassical.  The exactness for $\Sbeta$ is \statusproved.}
\end{statusbox}

\section{The complete finite vanishing spectrum}
\label{sec:vanishingspectrum}

Define
\[
 \Vclass_\beta=
 \left\{\tau_2(f):
 f\not\equiv0,\ f|_{\Sbeta}=0,\ \tau_2(f)<\infty\right\}.
\]
The sharp threshold identifies the left endpoint.  A simple rotated Gaussian factor fills
the whole remaining half-line.

\begin{theorem}[Exact vanishing spectrum]
\label{thm:vanishingspectrum}
One has
\begin{equation}
 \Vclass_\beta=[\Cbeta,\infty).
 \label{eq:vanishingspectrum}
\end{equation}
More explicitly, for every $a\ge0$ the function
\begin{equation}
 H_{\beta,a}(z)=\exp(-ia z^2)H_\beta(z)
 \label{eq:Hba}
\end{equation}
vanishes on $\Sbeta$ and satisfies
\[
 \tau_2(H_{\beta,a})=\Cbeta+a.
\]
\end{theorem}

\begin{proof}
The lower inclusion follows from Theorem~\ref{thm:sharpthreshold}.  For the reverse
inclusion, the upper bound
\[
 M(H_{\beta,a},r)\le \ee^{a r^2}M(H_\beta,r)
\]
gives $\tau_2(H_{\beta,a})\le a+\Cbeta$.  On the ray
$z=r\ee^{i\pi/4}$ one has $z^4=-r^4$, so Lemma~\ref{lem:maxmod} gives
$|H_\beta(z)|=M(H_\beta,r)$.  Also $z^2=ir^2$, hence
\[
 |\exp(-ia z^2)|=\ee^{a r^2}.
\]
Therefore equality holds in the type estimate.
\end{proof}

\begin{corollary}[No spectral gap above the zero-set threshold]
Finite positive types of vanishing functions are not exceptional: every type at or above
$\pi/(2\beta)$ occurs.  What is exceptional is the entire open interval below that value,
where the only vanishing function is zero.
\end{corollary}

\begin{remark}[No appeal to complete regularity is needed]
Recent angular-density theory predicts the same critical value and describes extremal
indicators geometrically.  The special fourfold product here is simpler: the exact
maximum-modulus ray and the rotated Gaussian prove the full spectrum without invoking a
classification of completely regular growth.
\end{remark}

\section{The corrected integer-valued type spectrum and the Pila gap}
\label{sec:pilagap}

\subsection{The spectrum of integer-valued entire functions}

Let
\[
 \Tclass_\beta=
 \left\{\tau_2(g):
 g\text{ entire},\ g(\Sbeta)\subseteq\Z,\ \tau_2(g)<\infty\right\}.
\]
The predecessor report suggested that natural constructions might realize only type zero
or infinite type.  The rotational products show otherwise.

\begin{theorem}[A large explicit part of the integer-valued spectrum]
\label{thm:integerspectrum}
One has
\begin{equation}
 \{0\}\cup[\Cbeta,\infty)\subseteq\Tclass_\beta.
 \label{eq:integerspectrum}
\end{equation}
For $\tau\ge\Cbeta$, an explicit representative is
\begin{equation}
 g_\tau(z)=2^z+
 \exp\bigl(-i(\tau-\Cbeta)z^2\bigr)H_\beta(z).
 \label{eq:gtau}
\end{equation}
\end{theorem}

\begin{proof}
Type zero is represented by constants, polynomials in $2^z$ and $3^z$, and many other
order-one functions.  For $s\in\Sbeta$, $H_\beta(s)=0$ and $2^s\in\Z$, so
$g_\tau(s)\in\Z$.  The second summand has type $\tau$ by
Theorem~\ref{thm:vanishingspectrum}; the first has order one.  Along the maximum ray of the
second summand its magnitude is $\exp(\tau r^2+o(r^2))$, while $|2^z|=\exp(O(r))$, so no
leading cancellation is possible.  Hence $\tau_2(g_\tau)=\tau$.
\end{proof}

\subsection{Pila's sufficient threshold}

Pila's theorem on entire functions sharing arguments of integrality specializes to the
following conditional statement.  It is a classical input, not reproved here.

\begin{theorem}[Pila, specialized]
\label{thm:pila}
Assume $\beta$ is a counterexample generator.  If an entire function $g$ is integer-valued
on $\Sbeta$ and
\begin{equation}
 0<\tau_2(g)\le
 \TP:=\frac{1}{832\beta^2\log2\log3},
 \label{eq:TP}
\end{equation}
then $2^z,3^z,g(z)$ are algebraically dependent.
\end{theorem}

The predecessor report also proves the following growth consequence.

\begin{proposition}[Algebraic dependence forces quadratic type zero]
\label{prop:algdepzero}
If an entire $g$ is algebraic over $\C(2^z,3^z)$, then $g$ has order at most one in the
Nevanlinna sense and in particular $\tau_2(g)=0$.
\end{proposition}

Combining the two gives the contradiction sought by the Pila programme: any construction
with \eqref{eq:TP} proves the Alaoglu--Erd\H{o}s conjecture.

\begin{corollary}[The true unresolved type interval]
\label{cor:trueinterval}
Conditional on a counterexample,
\begin{equation}
 \Tclass_\beta\cap(0,\TP]=\varnothing,
 \qquad
 [\Cbeta,\infty)\subseteq\Tclass_\beta.
 \label{eq:trueinterval}
\end{equation}
Thus the only unknown finite positive types lie in
\begin{equation}
 \boxed{
 \left(\frac{1}{832\beta^2\log2\log3},
       \frac{\pi}{2\beta}\right).}
 \label{eq:middleinterval}
\end{equation}
\end{corollary}

\subsection{The exact size of the gap}

The ratio between the least vanishing type and Pila's threshold is
\begin{equation}
 \frac{\Cbeta}{\TP}
 =416\pi\beta\log2\log3
 =995.206300791\ldots\,\beta.
 \label{eq:gapratio}
\end{equation}
The reciprocal is
\[
 \frac{\TP}{\Cbeta}
 =\frac1{416\pi\beta\log2\log3}.
\]

\begin{center}
\begin{tabular}{@{}rrrrr@{}}
\toprule
$\beta$ & $\Cbeta$ & $\TP$ & $\Cbeta/\TP$ & $\TP/\Cbeta$\\
\midrule
$14$ & $1.1219973763\times10^{-1}$ & $8.0528700100\times10^{-6}$
     & $1.3932888211\times10^4$ & $7.1772627818\times10^{-5}$\\
$27$ & $5.8177641733\times10^{-2}$ & $2.1651063401\times10^{-6}$
     & $2.6870570121\times10^4$ & $3.7215436646\times10^{-5}$\\
\bottomrule
\end{tabular}
\end{center}

\begin{remark}[Comparison with the earlier Jensen estimate]
The predecessor's elementary Jensen lower bound for a vanishing function was
$1/(4\beta)$.  The exact value $\pi/(2\beta)$ is larger by the factor $2\pi$.  More
importantly, the earlier claim of an infinite separation is replaced by the finite and exact
ratio \eqref{eq:gapratio}.
\end{remark}

\subsection{A decisive methodological consequence}

\begin{corollary}[Null-function correction cannot enter the Pila regime]
\label{cor:nullcorrection}
Let $F$ be any entire function and let $H\not\equiv0$ vanish on $\Sbeta$.  If both have
finite quadratic type, then $\tau_2(H)\ge\Cbeta$.  In particular, adding a nonzero null
function to an interpolant cannot preserve any type bound below $\Cbeta$, let alone Pila's
threshold $\TP$.
\end{corollary}

This closes a broad class of strategies: one cannot first build a convenient small function
with approximately integral values and then repair all values by adding a low-growth entire
function vanishing on the grid.  Below $\Cbeta$, interpolation data determine the function
uniquely.

\begin{statusbox}{The spectrum inclusion and exact gap are \statusproved.  The exclusion
$(0,\TP]$ uses Pila's theorem \statusclassical and the inherited algebraic-growth lemma.}
\end{statusbox}

\section{Carleman density and finite zero covers}
\label{sec:cover}

The sharp threshold has a local form that measures subsets of the monoid.  It is the right
tool for finite templates, alphabets, and decoder complexity.

\begin{definition}[Quadratic Carleman density]
For a discrete set $A\subset(1,\infty)$, define
\begin{equation}
 \DC(A)=\limsup_{R\to\infty}\frac1{\log R}
 \sum_{\substack{a\in A\\a^2<R}}
 \left(\frac1{a^2}-\frac{a^2}{R^2}\right).
 \label{eq:DC}
\end{equation}
\end{definition}

For a set with $N_A(r)=d r^2+O(r)$, the calculation in
Theorem~\ref{thm:raydensity} gives $\DC(A)=d$.  In particular,
\begin{equation}
 \DC(\Sbeta\setminus\{0,1\})=d_\beta=\frac1{2\beta}.
 \label{eq:DCS}
\end{equation}
Removing or adding finitely many points does not change the density.

\begin{theorem}[Zero-density bound]
\label{thm:DCzero}
If $h\not\equiv0$ is entire of finite quadratic type $\tau$ and $A\subset(1,\infty)$ is a
subset of its positive real zeros, then
\begin{equation}
 \DC(A)\le\frac{\tau}{\pi}.
 \label{eq:DCzero}
\end{equation}
\end{theorem}

\begin{proof}
Repeat the proof of Theorem~\ref{thm:raydensity} without assuming an asymptotic counting
law.  Lemma~\ref{lem:carleman} directly gives \eqref{eq:DCzero}.
\end{proof}

\begin{theorem}[Finite zero-cover inequality]
\label{thm:zerocover}
Let $h_1,\dots,h_K$ be nonzero entire functions of finite quadratic types
$\tau_1,\dots,\tau_K$.  If
\[
 \Sbeta\subseteq\bigcup_{j=1}^{K}\{s\in\R:h_j(s)=0\},
\]
then
\begin{equation}
 \boxed{\Cbeta\le\sum_{j=1}^{K}\tau_j.}
 \label{eq:coverineq}
\end{equation}
\end{theorem}

\begin{proof}
At every finite $R$, the nonnegative Carleman weight contributed by $\Sbeta$ is at most the
sum of the corresponding weights over the positive zeros of the $h_j$.  Divide by
$\log R$, take limsups, use Theorem~\ref{thm:DCzero}, and multiply by $\pi$.
\end{proof}

\begin{corollary}[Multiplicity form]
If every $s\in\Sbeta$ is a zero of total multiplicity at least $q$ across the functions
$h_1,\dots,h_K$, then
\[
 q\Cbeta\le\sum_{j=1}^{K}\tau_j.
\]
In particular, a single function vanishing to multiplicity at least $q$ on $\Sbeta$ has
type at least $q\Cbeta$.
\end{corollary}

\subsection{Template and coincidence bounds}

\begin{theorem}[Finite-template bound]
\label{thm:templates}
Let $g,F_1,\dots,F_K$ be entire functions of finite quadratic type.  Assume
$g\not\equiv F_j$ for every $j$ and that, for every $s\in\Sbeta$, one has
\[
 g(s)\in\{F_1(s),\dots,F_K(s)\}.
\]
Then
\begin{equation}
 \Cbeta\le
 \sum_{j=1}^{K}\max\{\tau_2(g),\tau_2(F_j)\}.
 \label{eq:templatebound}
\end{equation}
If all $F_j$ have type zero and $\tau_2(g)=\tau>0$, then
\begin{equation}
 K\tau\ge\Cbeta.
 \label{eq:Ktau}
\end{equation}
\end{theorem}

\begin{proof}
The nonzero residuals $h_j=g-F_j$ have type at most the displayed maxima, and their zero
sets cover $\Sbeta$.  Apply Theorem~\ref{thm:zerocover}.
\end{proof}

\begin{corollary}[Coincidence density with one model]
\label{cor:coincidence}
Let $F$ have type zero and let $g\not\equiv F$ have type $\tau<\infty$.  Then the
coincidence set
\[
 A_F=\{s\in\Sbeta:g(s)=F(s)\}
\]
satisfies
\[
 \DC(A_F)\le\frac{\tau}{\pi}.
\]
Relative to the full monoid Carleman density, the proportion is at most
\[
 \frac{\DC(A_F)}{\DC(\Sbeta)}\le\frac{\tau}{\Cbeta}.
\]
\end{corollary}

At Pila's threshold this relative upper bound is
\[
 \frac{\TP}{\Cbeta}
 =\frac1{416\pi\beta\log2\log3}.
\]
For $\beta=14$ it is approximately $7.18\times10^{-5}$, or $0.00718\%$.
Thus a Pila-killing function cannot be a sparse perturbation of any fixed type-zero model:
it must disagree with that model on more than $99.9928\%$ of the monoid in Carleman density.

\begin{statusbox}{The zero-density, cover, template, and coincidence inequalities are
\statusproved.}
\end{statusbox}

\section{Finite alphabets, carries, and local decoder complexity}
\label{sec:alphabet}

The cover theorem is especially effective when the covering functions are obtained by
subtracting finitely many possible values from one type-preserving observable.

\subsection{A general alphabet lemma}

For $h\in\C$, let $\Tsh{h}$ denote translation:
\[
 (\Tsh{h}g)(z)=g(z+h).
\]
A finite translation operator is an expression
\[
 L=\sum_{\nu=1}^{r}c_\nu\Tsh{h_\nu},
 \qquad c_\nu,h_\nu\in\C.
\]
Translations and finite linear combinations do not increase quadratic type.

\begin{theorem}[Finite-alphabet dichotomy]
\label{thm:alphabet}
Let $g$ be entire with finite quadratic type $\tau$, let $L$ be a finite translation
operator, and let $A\subset\C$ be finite.  Suppose
\[
 (Lg)(s)\in A\qquad(s\in\Sbeta).
\]
Then either
\begin{enumerate}
\item $Lg\equiv a$ for some $a\in A$, or
\item $|A|\tau\ge\Cbeta$.
\end{enumerate}
\end{theorem}

\begin{proof}
For $a\in A$, put $h_a=Lg-a$.  The zero sets of the $h_a$ cover $\Sbeta$.  Each nonzero
$h_a$ has type at most $\tau$.  If none is identically zero, apply
Theorem~\ref{thm:zerocover}.
\end{proof}

\begin{corollary}[Value alphabet]
\label{cor:valuealphabet}
Let $g$ be a nonconstant entire function of positive finite quadratic type $\tau$ and
suppose $g(\Sbeta)\subseteq A$ for a finite set $A$.  Then
\[
 |A|\ge\frac{\Cbeta}{\tau}.
\]
If $g$ is integer-valued, the same statement applies to any finite integer value alphabet.
\end{corollary}

\begin{corollary}[Pila-scale state explosion]
\label{cor:stateexplosion}
If $0<\tau\le\TP$, every finite value alphabet for a nonconstant integer-valued $g$ has
cardinality at least
\begin{equation}
 \left\lceil\frac{\Cbeta}{\tau}\right\rceil
 \ge
 \left\lceil416\pi\beta\log2\log3\right\rceil.
 \label{eq:states}
\end{equation}
At $\beta=14$, the right side is $13{,}933$; at $\beta=27$, it is $26{,}871$.
\end{corollary}

\subsection{Coordinate increments}

For the sample array
\[
 u_{n,k}=g(n+k\beta),
\]
the two coordinate increments are
\[
 \Del{1}g=\Tsh{1}g-g,
 \qquad
 \Del{\beta}g=\Tsh{\beta}g-g.
\]
If $g$ is integer-valued on $\Sbeta$, both increments are integer-valued there.

\begin{corollary}[Increment-alphabet dichotomy]
\label{cor:incrementalphabet}
Suppose $\Del{1}g(\Sbeta)$ is contained in a finite set $A_1$ and
$\Del{\beta}g(\Sbeta)$ is contained in a finite set $A_\beta$.  Then, for each
$h\in\{1,\beta\}$, either
\[
 \Del{h}g\equiv a_h\quad\text{for some }a_h\in A_h,
\]
or
\[
 |A_h|\tau_2(g)\ge\Cbeta.
\]
\end{corollary}

\begin{proposition}[Two globally constant increments force constancy]
\label{prop:twoconstantincrements}
Let $g$ be entire and integer-valued on $\Sbeta$.  If
\[
 \Del{1}g\equiv a,
 \qquad
 \Del{\beta}g\equiv b
\]
for constants $a,b$, then $g$ is a constant integer.
\end{proposition}

\begin{proof}
The values of the increments on $\Sbeta$ show $a,b\in\Z$.  Put $h(z)=g(z)-az$.
Then $h$ is $1$-periodic and
\[
 h(z+\beta)-h(z)=b-a\beta=:c.
\]
On the real line, $h$ is bounded because it is continuous and $1$-periodic.  Iterating the
last identity gives $h(x+k\beta)=h(x)+kc$, so $c=0$.  Thus $h$ has the two real periods
$1$ and $\beta$.  Their integer span is dense in $\R$, hence continuity makes $h$ constant
on every horizontal line and the identity theorem makes it constant on $\C$.  Finally,
$b=a\beta$ with $a,b\in\Z$ and irrational $\beta$ forces $a=b=0$.
\end{proof}

\begin{corollary}[At least one coordinate has a large alphabet]
\label{cor:oneaxislarge}
Let $g$ be integer-valued on $\Sbeta$ with $0<\tau_2(g)=\tau<\Cbeta$.  If both coordinate
increment alphabets are finite, then at least one has cardinality at least
$\lceil\Cbeta/\tau\rceil$.
\end{corollary}

\begin{proof}
If both alphabets had cardinality strictly below $\Cbeta/\tau$, Corollary
~\ref{cor:incrementalphabet} would make both increments globally constant, contradicting
Proposition~\ref{prop:twoconstantincrements} and positive type.
\end{proof}

\subsection{Interpretation for carry languages}

Many attempted constructions encode $g(n+k\beta)$ by a finite-state carry process.  A local
observable of such a process typically has a finite alphabet after applying a fixed
translation operator: a first difference, a short stencil, a carry defect, or a residue
correction.  Theorem~\ref{thm:alphabet} gives a universal test.

\begin{itemize}
\item If the alphabet is smaller than $\Cbeta/\tau$, the observable is not merely simple
on the monoid: it is a global entire identity.
\item If the global identity is incompatible with irrationality or integrality, the
construction is impossible.
\item At Pila type, a surviving finite-state observable needs on the order of
$10^4$ states even before any arithmetic height or denominator cost is imposed.
\end{itemize}

This is stronger than a Shannon-entropy heuristic.  It is an exact zero-cover obstruction
coming from the angular concentration of the interpolation nodes.

\begin{statusbox}{All alphabet and increment statements are \statusproved.}
\end{statusbox}

\section{Polynomial coordinate decoders are rigid}
\label{sec:polydecoder}

A natural attempt is to prescribe the integer values of $g$ by a low-complexity formula in
the unique coordinates $(n,k)$ of $s=n+k\beta$.  Below the exact zero-set threshold, this
is impossible except when the formula already comes from a one-variable polynomial in
$s$.

\subsection{Two finite-difference equations}

Write
\[
 \Del{h}f(z)=f(z+h)-f(z),
 \qquad
 \Del{h}^{r}=\underbrace{\Del{h}\circ\cdots\circ\Del{h}}_{r\text{ times}}.
\]

\begin{lemma}[Two incommensurable polynomial-difference equations]
\label{lem:twodifference}
Let $\beta\notin\Q$.  If an entire function $G$ satisfies
\[
 \Del{1}^{p+1}G\equiv0,
 \qquad
 \Del{\beta}^{q+1}G\equiv0
\]
for nonnegative integers $p,q$, then $G$ is a polynomial of degree at most
$\min\{p,q\}$.
\end{lemma}

\begin{proof}
The equation $\Del{1}^{p+1}G=0$ has the standard Floquet--Newton form
\begin{equation}
 G(z)=\sum_{j=0}^{p}z^j P_j(z),
 \label{eq:floquetpoly}
\end{equation}
where every $P_j$ is entire and $1$-periodic.  This follows by induction: the $p$th first
difference is periodic, subtract a suitable binomial polynomial times that periodic
function, and descend.

Each $P_j$ has a locally normally convergent Fourier--Laurent expansion
\[
 P_j(z)=\sum_{m\in\Z}c_{j,m}\ee^{2\pi i m z}.
\]
Collecting powers of $z$ gives
\[
 G(z)=\sum_{m\in\Z}\ee^{2\pi i m z}Q_m(z),
 \qquad \deg Q_m\le p.
\]
On the $m$th mode, $\Del{\beta}$ acts on polynomials by
\[
 Q(z)\longmapsto
 \ee^{2\pi i m\beta}Q(z+\beta)-Q(z).
\]
For $m\ne0$, irrationality gives $\ee^{2\pi i m\beta}\ne1$.  The last operator is a
nonzero scalar multiple of the identity plus a nilpotent operator on the finite-dimensional
space of polynomials of degree at most $p$; it is therefore invertible.  The equation
$\Del{\beta}^{q+1}G=0$ forces every $Q_m$ with $m\ne0$ to vanish.  The remaining mode
$Q_0$ is a polynomial, and $\Del{\beta}^{q+1}Q_0=0$ gives $\deg Q_0\le q$.
\end{proof}

\subsection{The decoder theorem}

\begin{theorem}[Polynomial decoder rigidity]
\label{thm:polydecoder}
Let $G$ be entire with
\[
 \tau_2(G)<\Cbeta.
\]
Suppose there is a polynomial $P(X,Y)\in\C[X,Y]$ such that
\begin{equation}
 G(n+k\beta)=P(n,k)
 \qquad(n,k\in\Nzero).
 \label{eq:decoder}
\end{equation}
Then there is a one-variable polynomial $Q\in\C[Z]$ for which
\begin{equation}
 G(z)=Q(z),
 \qquad
 P(X,Y)=Q(X+\beta Y).
 \label{eq:factorline}
\end{equation}
\end{theorem}

\begin{proof}
Let $p=\deg_XP$ and $q=\deg_YP$.  Evaluating the $(p+1)$st first difference at
$s=n+k\beta$ gives
\[
 (\Del{1}^{p+1}G)(s)=0
 \qquad(s\in\Sbeta).
\]
This entire residual has type at most $\tau_2(G)<\Cbeta$, so subcritical uniqueness makes
it identically zero.  Similarly $\Del{\beta}^{q+1}G\equiv0$.  Lemma
~\ref{lem:twodifference} gives $G=Q$ for a polynomial $Q$.

The polynomial
\[
 R(X,Y)=P(X,Y)-Q(X+\beta Y)
\]
vanishes on $\Nzero^2$.  Fixing $Y=k$ shows $R(X,k)$ has infinitely many zeros in $X$,
so it is zero.  Its coefficients as polynomials in $X$ then vanish at infinitely many $k$,
therefore $R\equiv0$.
\end{proof}

\begin{corollary}[Algebraic-coefficient decoders are constant]
\label{cor:algdecoder}
Assume, as in the conditional structure theorem, that $\beta$ is transcendental.  In
Theorem~\ref{thm:polydecoder}, if $P\in\Abar[X,Y]$, then $P$ and $G$ are constant.  If
$G(\Sbeta)\subseteq\Z$, the constant is an integer.
\end{corollary}

\begin{proof}
Suppose $Q$ in \eqref{eq:factorline} has degree $d\ge1$ and leading coefficient $c\ne0$.
Then the coefficient of $X^d$ in $P$ is $c$, while the coefficient of $X^{d-1}Y$ is
$dc\beta$.  Their ratio is $d\beta$, which cannot be algebraic.  Hence $d=0$.
\end{proof}

\begin{corollary}[No rational affine coordinate code]
A subcritical entire function cannot satisfy
\[
 G(n+k\beta)=an+bk+c
\]
with $a,b,c\in\Q$ unless $a=b=0$ and $G\equiv c$.
\end{corollary}

\begin{remark}
The theorem is not a statement that every low-type integer-valued function is polynomial.
It says that if its values admit a polynomial decoder in the hidden coordinates $(n,k)$,
then the decoder must forget those coordinates and depend only on the actual complex
argument $n+k\beta$.  Transcendence of $\beta$ then kills every algebraic-coefficient
nonconstant decoder.
\end{remark}

\begin{statusbox}{The finite-difference lemma and polynomial decoder theorem are
\statusproved.}
\end{statusbox}

\section{Constant-coefficient recurrences and Hankel-rank collapse}
\label{sec:recurrence}

A second common low-complexity model is that each row and each column of the sample array
satisfies a fixed linear recurrence.  In the entire-function setting, two incommensurable
recurrences leave only exponential polynomials.

\subsection{From sample recurrences to global equations}

Let
\[
 A(U)=\sum_{j=0}^{r}a_jU^j,
 \qquad
 B(V)=\sum_{\ell=0}^{s}b_\ell V^\ell
\]
be nonzero polynomials.  The array $u_{n,k}=G(n+k\beta)$ satisfies the corresponding
coordinate recurrences if
\begin{align}
 \sum_{j=0}^{r}a_j u_{n+j,k}&=0,
 \label{eq:nrec}\\
 \sum_{\ell=0}^{s}b_\ell u_{n,k+\ell}&=0
 \label{eq:krec}
\end{align}
for all $n,k\ge0$.  Equivalently, the residual entire functions
$A(\Tsh{1})G$ and $B(\Tsh{\beta})G$ vanish on $\Sbeta$.

\begin{lemma}[Two incommensurable linear difference equations]
\label{lem:twoCfinite}
Let $\beta\notin\Q$.  If an entire function $G$ satisfies
\[
 A(\Tsh{1})G\equiv0,
 \qquad
 B(\Tsh{\beta})G\equiv0
\]
for nonzero constant-coefficient polynomials $A,B$, then $G$ is an exponential polynomial:
\begin{equation}
 G(z)=\sum_{\nu=1}^{N}P_\nu(z)\ee^{\lambda_\nu z}
 \label{eq:exppoly}
\end{equation}
with polynomials $P_\nu$.  In particular, $G$ has order at most one and
$\tau_2(G)=0$.
\end{lemma}

\begin{proof}
Factor $A$.  Standard one-step difference theory writes the solutions of
$A(\Tsh{1})G=0$ as a finite sum of terms
\[
 \ee^{\lambda z}\sum_{j=0}^{m-1}z^jP_j(z),
\]
where $\ee^\lambda$ is a root of $A$, $m$ is its multiplicity, and each $P_j$ is entire and
$1$-periodic.  Expand the $P_j$ into Fourier modes $\ee^{2\pi i n z}$.

For a fixed root branch $\lambda$ and Fourier index $n$, the second shift has multiplier
\[
 \ee^{\beta(\lambda+2\pi i n)}.
\]
The equation $B(\Tsh{\beta})G=0$ permits only multipliers that are roots of $B$, with a
bounded polynomial multiplicity.  For fixed $\lambda$, the numbers
$\ee^{\beta(\lambda+2\pi i n)}$ are pairwise distinct because $\beta$ is irrational.
Since $B$ has finitely many roots, only finitely many Fourier indices survive.  This gives
\eqref{eq:exppoly}.
\end{proof}

\begin{theorem}[C-finite sample arrays have type zero]
\label{thm:Cfinite}
Let $G$ be entire, integer-valued on $\Sbeta$, and satisfy
\[
 \tau_2(G)<\Cbeta.
\]
If its sample array satisfies nontrivial constant-coefficient recurrences
\eqref{eq:nrec} and \eqref{eq:krec}, then $\tau_2(G)=0$.
\end{theorem}

\begin{proof}
The two residuals have type at most $\tau_2(G)$ and vanish on $\Sbeta$.  Subcritical
uniqueness makes both global equations.  Lemma~\ref{lem:twoCfinite} finishes the proof.
\end{proof}

\begin{corollary}[Positive subcritical type forces infinite recurrence complexity]
\label{cor:infiniteHankel}
If $0<\tau_2(G)<\Cbeta$, then the sample array cannot have finite constant-coefficient
recurrence rank in both coordinate directions.  Equivalently, at least one family of
one-dimensional coordinate sequences has unbounded Hankel rank.
\end{corollary}

\subsection{A nonlinear first-order no-go}

The same uniqueness mechanism globalizes sufficiently low-degree polynomial recursions.

\begin{proposition}[Polynomial iteration of degree at least two]
\label{prop:polyiteration}
Let $P\in\C[X]$ have degree $d\ge2$.  A nonconstant entire function of finite order cannot
satisfy
\[
 G(z+h)=P(G(z))
\]
for a nonzero $h\in\C$.
\end{proposition}

\begin{proof}
Choose $w_0=G(z_0)$ sufficiently large.  For large $|w|$, one has
$|P(w)|\ge c|w|^d$ with $c>0$.  Iterating along $z_0+nh$ gives
\[
 \log|G(z_0+nh)|\ge c_1d^n-c_2.
\]
Since $|z_0+nh|=O(n)$, the maximum modulus has infinite order, a contradiction.
\end{proof}

\begin{corollary}[Sample polynomial iteration]
\label{cor:samplepolyiteration}
Let $G$ be entire with finite quadratic type $\tau$, and suppose
\[
 G(s+h)=P(G(s))
\]
for every $s\in\Sbeta$, where $h\in\{1,\beta\}$ and $P$ has degree $d\ge2$.  If
$d\tau<\Cbeta$, then $G$ is constant.
\end{corollary}

\begin{proof}
The residual $G(z+h)-P(G(z))$ has type at most $d\tau$, vanishes on $\Sbeta$, and is
therefore identically zero.  Proposition~\ref{prop:polyiteration} forces constancy.
\end{proof}

\begin{remark}
Affine one-shift equations are not excluded: finite-order periodic functions can have
quadratic growth.  The second incommensurable recurrence is essential in
Theorem~\ref{thm:Cfinite}.  This accurately marks where the nonlinear difference-algebra
problem begins.
\end{remark}

\begin{statusbox}{The recurrence and polynomial-iteration results are \statusproved.}
\end{statusbox}

\section{Pila's multi-function theorem and a many-interpolant gate}
\label{sec:multipila}

The one-function Pila threshold is extremely small because one additional function must
supply all of the auxiliary-monomial volume.  Pila's general theorem allows several
additional integer-valued functions.  Specializing that theorem to $\Sbeta$ produces a
family of less severe, higher-dimensional targets that does not appear in the predecessor
report.

\subsection{Normalization on the ordered monoid}

Let
\[
 0=x_0<x_1<x_2<\cdots
\]
be the increasing enumeration of $\Sbeta$.  Proposition~\ref{prop:counting} gives
\begin{equation}
 x_n\sim\sqrt{2\beta n}.
 \label{eq:xnasymp}
\end{equation}
For an entire $g$ of finite quadratic type $\tau$, define the convenient upper normalized
rate
\begin{equation}
 \lambda_\beta(g):=2\beta\tau.
 \label{eq:lambdanorm}
\end{equation}
Then
\[
 \limsup_{n\to\infty}\frac{\log M(g,x_n)}n
 \le\lambda_\beta(g).
\]
For the two basic functions, Pila's scale-growth constants are
\[
 \omega_1=\sqrt{2\beta}\log2,
 \qquad
 \omega_2=\sqrt{2\beta}\log3,
 \qquad
 \omega_1\omega_2=2\beta\log2\log3.
\]

Pila's scale-separation constant obeys
\[
 \chi(\Sbeta,2)\ge\chi_2(2),
\]
where
\begin{equation}
 \chi_2(2)=
 1-\sqrt2+\log2-\log(\sqrt2-1)
 =1.1603072052\ldots.
 \label{eq:chi2}
\end{equation}
We fix the rational safety margin
\begin{equation}
 A_0:=\frac{29}{75}=0.386666\ldots
 <\frac{\chi_2(2)}3.
 \label{eq:A0}
\end{equation}

\subsection{The weighted-simplex criterion}

\begin{theorem}[Multi-function Pila gate]
\label{thm:multipila}
Assume a counterexample generator $\beta$ exists.  Let
$g_1,\dots,g_q$ be entire functions integer-valued on $\Sbeta$, with finite quadratic
types $\tau_1,\dots,\tau_q$, and put
\[
 \lambda_i=2\beta\tau_i.
\]
If
\begin{equation}
 \sum_{\substack{\mathbf j\in\Nzero^q\\
                  \mathbf j\cdot\boldsymbol\lambda\le A_0}}
 \bigl(A_0-\mathbf j\cdot\boldsymbol\lambda\bigr)^2
 >8\beta\log2\log3,
 \label{eq:multipilacriterion}
\end{equation}
then the entire functions
\[
 2^z,\ 3^z,\ g_1(z),\dots,g_q(z)
\]
are algebraically dependent over $\C$.
\end{theorem}

\begin{proof}
Apply Pila's Theorem~4.1 with two base functions, $t=2$, $s=2$, and $A=A_0$.  The
right side of Pila's monomial-count inequality is
\[
 s\,2!\,\omega_1\omega_2
 =2\cdot2\cdot2\beta\log2\log3
 =8\beta\log2\log3,
\]
which gives \eqref{eq:multipilacriterion}.  Pila's construction yields a nonzero integral
polynomial relation on the set $\Sbeta$, with every occurring monomial
$g_1^{j_1}\cdots g_q^{j_q}$ satisfying
$\mathbf j\cdot\boldsymbol\lambda\le A_0$.

After composition, the resulting entire residual has quadratic type at most
$A_0/(2\beta)$.  Since $A_0<\pi$, this is strictly smaller than
$\Cbeta=\pi/(2\beta)$.  The residual vanishes on $\Sbeta$, so
Corollary~\ref{cor:uniqueness} makes it identically zero.  Thus the dependence is an
identity of entire functions, not merely a relation on the interpolation set.
\end{proof}

\begin{remark}[Why the sharp zero threshold matters]
Pila's original proof already used a weaker Jensen uniqueness estimate to promote the
relation to an entire identity in the one-function case.  Theorem~\ref{thm:sharpthreshold}
shows conceptually that every weighted auxiliary polynomial produced with budget
$A_0<\pi$ lies safely inside the exact uniqueness region.  No separate case-by-case Jensen
calculation is needed for several additional functions.
\end{remark}

\subsection{Equal-type functions}

If all $g_i$ have the same type $\tau$, put
\[
 \lambda=2\beta\tau,
 \qquad
 L=\left\lfloor\frac{A_0}{\lambda}\right\rfloor.
\]
The left side of \eqref{eq:multipilacriterion} becomes the exact finite sum
\begin{equation}
 \Phi_q(A_0,\lambda)
 =\sum_{\ell=0}^{L}
   \binom{\ell+q-1}{q-1}(A_0-\ell\lambda)^2.
 \label{eq:Phiq}
\end{equation}
Indeed, the number of $q$-tuples of nonnegative integers with coordinate sum $\ell$ is
$\binom{\ell+q-1}{q-1}$.

\begin{corollary}[Exact equal-type gate]
\label{cor:equaltypegate}
If $g_1,\dots,g_q$ are algebraically independent together with $2^z,3^z$, then necessarily
\begin{equation}
 \Phi_q(A_0,2\beta\tau)
 \le8\beta\log2\log3.
 \label{eq:equaltypenecessary}
\end{equation}
Consequently, constructing algebraically independent functions for which the reverse strict
inequality holds proves the Alaoglu--Erd\H{o}s conjecture.
\end{corollary}

For fixed $q$ and $\lambda\downarrow0$, a simplex Riemann sum gives
\begin{equation}
 \Phi_q(A_0,\lambda)
 =\frac{2A_0^{q+2}}{(q+2)!\lambda^q}
  +O_{A_0,q}(\lambda^{1-q}).
 \label{eq:Phiqasymp}
\end{equation}
Thus the natural equal-type target scale is
\begin{equation}
 \tau
 \asymp_q
 \frac1{2\beta}
 \left(
  \frac{A_0^{q+2}}
       {4(q+2)!\,\beta\log2\log3}
 \right)^{1/q}
 \asymp_q\beta^{-1-1/q}.
 \label{eq:qscale}
\end{equation}
For $q=1$ this is the familiar $\beta^{-2}$ scale.  For two or more algebraically
independent functions, the power of $\beta$ improves.

\subsection{Numerical thresholds}

The following table solves the exact inequality \eqref{eq:Phiq} by bisection.  Each entry
is the supremal common quadratic type $\tau$ for which the strict Pila inequality holds;
a rigorous application uses any slightly smaller value.

\begin{center}
\begin{tabular}{@{}rcc@{}}
\toprule
$q$ & supremal $\tau$ at $\beta=14$ & supremal $\tau$ at $\beta=27$\\
\midrule
$1$ & $8.07651\times10^{-6}$ & $2.17055\times10^{-6}$\\
$2$ & $1.71057\times10^{-4}$ & $6.34265\times10^{-5}$\\
$3$ & $4.61028\times10^{-4}$ & $1.88916\times10^{-4}$\\
$4$ & $7.55746\times10^{-4}$ & $3.24386\times10^{-4}$\\
$5$ & $1.02034\times10^{-3}$ & $4.49547\times10^{-4}$\\
$8$ & $1.63138\times10^{-3}$ & $7.45529\times10^{-4}$\\
\bottomrule
\end{tabular}
\end{center}

The $q=1$ entry is close to, and slightly improves on, the conservative
$1/832$ constant because \eqref{eq:Phiq} uses the exact finite monomial sum rather than an
integral lower bound.  The gain from several functions is substantial: at $\beta=14$, two
independent functions may each be about twenty-one times larger than the one-function Pila
threshold; eight may each be about two hundred times larger.

\begin{statusbox}{The specialization and exact finite-sum criterion use Pila's theorem
\statusclassical; the normalization, promotion to an entire identity, asymptotic, and
numerical evaluation are \statusproved and \statuscomputed.}
\end{statusbox}

\section{Translation orbits and finite transformal rank}
\label{sec:translationrank}

The multi-function gate becomes especially natural when all additional functions are
translates of one integer-valued interpolant.

\subsection{The translate family}

For $h\in\Sbeta$, define
\[
 g_h(z)=g(z+h).
\]
Because $\Sbeta+h\subseteq\Sbeta$, every $g_h$ is integer-valued on $\Sbeta$.  Translation
preserves quadratic type.  Moreover, the base field
\begin{equation}
 K_0=\C(2^z,3^z)
 \label{eq:K0}
\end{equation}
is stable under all shifts by $h\in\Sbeta$: the multipliers $2^h$ and $3^h$ are integers.

\begin{theorem}[Quantitative algebraic translation-rank bound]
\label{thm:translationrank}
Let $g$ be entire and integer-valued on $\Sbeta$, with positive finite quadratic type $\tau$, and
put
\[
 \lambda=2\beta\tau.
\]
Assume $\lambda<A_0$.  Let $q_*(\beta,\lambda)$ be the least positive integer $q$ such that
\begin{equation}
 \Phi_q(A_0,\lambda)>8\beta\log2\log3.
 \label{eq:qstar}
\end{equation}
Then
\begin{equation}
 \trdeg_{K_0}
 K_0\bigl(g(z+h):h\in\Sbeta\bigr)
 \le q_*(\beta,\lambda)-1.
 \label{eq:trdegbound}
\end{equation}
\end{theorem}

\begin{proof}
If the relative transcendence degree were at least $q_*$, one could choose
$h_1,\dots,h_{q_*}\in\Sbeta$ so that the translates $g_{h_i}$ are algebraically
independent over $K_0$.  Then $2^z,3^z,g_{h_1},\dots,g_{h_{q_*}}$ would be algebraically
independent over $\C$, contradicting Theorem~\ref{thm:multipila} and
\eqref{eq:qstar}.
\end{proof}

A simpler explicit, though often weaker, bound uses only the zero vector and the $q$ unit
vectors in the weighted simplex.

\begin{corollary}[Unit-vector translation bound]
\label{cor:unittranslation}
If $\lambda<A_0$ and
\begin{equation}
 q>
 \frac{8\beta\log2\log3-A_0^2}{(A_0-\lambda)^2},
 \label{eq:unitq}
\end{equation}
then every $q$ translates of $g$ are algebraically dependent over $K_0$.
\end{corollary}

\begin{proof}
The sum \eqref{eq:Phiq} contains the contribution
$A_0^2+q(A_0-\lambda)^2$ from total degrees zero and one.
\end{proof}

\subsection{Exact sample bounds}

The exact $q_*$ from \eqref{eq:qstar} can be much smaller than the unit-vector estimate
when $A_0/\lambda$ permits higher total degrees.

\begin{center}
\begin{tabular}{@{}crrrr@{}}
\toprule
& \multicolumn{4}{c}{$\lambda/A_0$}\\
\cmidrule(l){2-5}
$\beta$ & $0.05$ & $0.10$ & $0.25$ & $0.50$\\
\midrule
$14$ & $4$ & $7$ & $34$ & $2278$\\
$27$ & $5$ & $8$ & $43$ & $4397$\\
\bottomrule
\end{tabular}
\end{center}

For example, at $\beta=14$ and normalized rate $\lambda=0.1A_0$, seven suitably chosen
translates cannot remain algebraically independent over $K_0$.

\subsection{The difference-algebraic consequence}

Let $\sigma_1$ and $\sigma_\beta$ be the commuting automorphisms induced by translation by
$1$ and $\beta$.  For type zero, Pila's theorem may be applied with arbitrarily small positive upper rates and already gives algebraicity over $K_0$; hence the genuinely new case is $\tau>0$.

The field generated by all transforms
\[
 \sigma_1^n\sigma_\beta^k(g)=g(z+n+k\beta)
\]
has finite transcendence degree over the shift-stable base field $K_0$ whenever
$2\beta\tau<A_0$.  In the language of difference algebra, $g$ is transformally algebraic
for the two-shift action.

\begin{corollary}[Existence of an algebraic partial-difference equation]
\label{cor:algdifference}
Under the hypotheses of Theorem~\ref{thm:translationrank}, there is a nonzero polynomial
$P$ over $K_0$ and finitely many pairs $(n_j,k_j)\in\Nzero^2$ such that
\begin{equation}
 P\bigl(g(z+n_1+k_1\beta),\dots,
       g(z+n_r+k_r\beta)\bigr)\equiv0.
 \label{eq:partialdifference}
\end{equation}
The number of transforms needed may be bounded in terms of
$q_*(\beta,2\beta\tau)$.
\end{corollary}

\begin{proof}
Finite relative transcendence degree means that any sufficiently large finite set of
transforms is algebraically dependent over $K_0$.
\end{proof}

\begin{target}[Two-shift finite-order classification]
\label{target:differenceclassification}
Prove that an entire function $g$ satisfying all of the following conditions must be
algebraic over $K_0=\C(2^z,3^z)$:
\begin{enumerate}
\item $g(\Sbeta)\subseteq\Z$;
\item $0<2\beta\tau_2(g)<A_0$;
\item the translation field
$K_0(g(z+n+k\beta):n,k\in\Nzero)$ has finite relative transcendence degree.
\end{enumerate}
Such a theorem, combined with Proposition~\ref{prop:algdepzero}, would prove the
Alaoglu--Erd\H{o}s conjecture.
\end{target}

The linear portion of this target is already closed by
Theorem~\ref{thm:Cfinite}; the polynomial coordinate-decoder portion is closed by
Theorem~\ref{thm:polydecoder}.  What remains is a genuinely nonlinear algebraic
partial-difference equation with two incommensurable real shifts.

\subsection{Why this route is not a disguised null-function construction}

The translate field uses channel information unavailable to $H_\beta$.  A null function
has type at least $\Cbeta$, so its normalized rate is at least
\[
 2\beta\Cbeta=\pi>A_0.
\]
It lies completely outside Theorem~\ref{thm:translationrank}.  The new target concerns a
direct integer-valued interpolant of normalized type less than $0.3867$, whereas the least
nonzero null function has normalized type $\pi\approx3.1416$.  The separation is intrinsic,
not a weak numerical constant.

\begin{statusbox}{The translation-rank theorem and partial-difference consequence are
\statusproved from the multi-function Pila gate.  Target~\ref{target:differenceclassification}
is \statusopen.}
\end{statusbox}

\section{Closed subcases of the two-shift programme}
\label{sec:closedsubcases}

The open difference-algebra target is broad, but several natural degeneracies are already
ruled out by elementary arguments.  These results help define the minimal genuinely
nonlinear case.

\subsection{Periodic branches}

\begin{proposition}[One real period plus integrality is fatal]
\label{prop:periodicfatal}
If an entire function $g$ is $1$-periodic and integer-valued on $\Sbeta$, then $g$ is a
constant integer.  The same conclusion holds if $g$ is $\beta$-periodic.
\end{proposition}

\begin{proof}
If $g$ is $1$-periodic, then $g(k\beta)\in\Z$ for every $k\ge0$.  The set
$\{k\beta\bmod1:k\ge0\}$ is dense in $[0,1]$.  By continuity, $g(x)$ is integer-valued
for every real $x$.  A continuous map from a connected interval to the discrete set $\Z$
is constant, and the identity theorem extends the constancy to $\C$.  The
$\beta$-periodic case uses the dense set $\{n\bmod\beta:n\ge0\}$.
\end{proof}

\begin{corollary}[Finite-dimensional translation span]
\label{cor:finitedimspan}
If
\[
 \dim_\C\Span\{g(z+n+k\beta):n,k\in\Nzero\}<\infty
\]
and $g$ is entire, integer-valued on $\Sbeta$, and subcritical, then
$\tau_2(g)=0$.
\end{corollary}

\begin{proof}
The two commuting translation operators act linearly on a finite-dimensional vector space,
so each satisfies a nonzero characteristic polynomial.  Apply
Theorem~\ref{thm:Cfinite}.
\end{proof}

\subsection{Algebraic rank zero}

\begin{proposition}[Relative algebraic rank zero]
\label{prop:rankzero}
If $g$ is entire and algebraic over $K_0=\C(2^z,3^z)$, then $\tau_2(g)=0$.
Consequently a positive-type interpolant in the translation-rank regime must have positive
relative transcendence degree.
\end{proposition}

\begin{proof}
This is Proposition~\ref{prop:algdepzero}.
\end{proof}

\subsection{Degree growth as the next obstruction}

Suppose a finite set of transforms of $g$ gives a meromorphic map
\[
 F:\C\longrightarrow X
\]
to an algebraic variety over $K_0$, and the two shifts descend to commuting dominant
rational self-maps $\Phi_1,\Phi_\beta$ of $X$:
\[
 F(z+1)=\Phi_1(F(z)),
 \qquad
 F(z+\beta)=\Phi_\beta(F(z)).
\]
For a scalar polynomial map of degree at least two,
Proposition~\ref{prop:polyiteration} shows that iteration forces infinite order.  The
natural higher-dimensional replacement is dynamical degree.

\begin{target}[Dynamical-degree reduction]
\label{target:dynamicaldegree}
Prove that if $F$ has finite order and Zariski-dense image, then both
$\Phi_1$ and $\Phi_\beta$ have first dynamical degree one.  More generally, show that any
component with dynamical degree greater than one is contained in a proper common exceptional
subvariety.
\end{target}

If this reduction holds, the surviving maps are birationally low-entropy.  The expected
models are translations or affine endomorphisms on algebraic groups, fibrations by such
models, and finite quotients.  The dense real subgroup generated by $1$ and $\beta$, together
with integer values on its positive cone, should then eliminate periodic and elliptic
branches and force algebraicity over $K_0$.

\begin{target}[Low-entropy two-shift classification]
\label{target:lowentropy}
Classify finite-order entire curves carrying two commuting, incommensurable real-shift
actions of dynamical degree one.  In the integer-valued setting, prove that every such curve
is generated by exponential polynomials and functions algebraic over $K_0$.
\end{target}

\subsection{Why one shift is not enough}

A single affine difference equation can support positive quadratic type.  For example, a
$1$-periodic entire function may have a rapidly decaying Fourier series
\[
 P(z)=\sum_{n\in\Z}c_n\ee^{2\pi i n z}
\]
whose maximum modulus grows like $\exp(O(r^2))$.  Thus finite order alone does not make one
shift finite-dimensional.  The second irrational shift is the rigidity source: it samples
the full Fourier spectrum through the dense multipliers $\ee^{2\pi i n\beta}$.

This observation explains why Theorem~\ref{thm:Cfinite} is robust and why the nonlinear
problem should be formulated for two commuting shifts from the outset.

\section{A general integer-order ray theorem}
\label{sec:generalray}

The fourfold product and the constant $\pi/(2\beta)$ are instances of a general principle
for positive-ray sets of integer exponent of convergence.  The proof is the same square-root
idea with a $k$th root.

\begin{definition}
For an entire function $f$ and an integer $k\ge1$, define its order-$k$ type by
\[
 \tau_k(f)=\limsup_{r\to\infty}\frac{\log M(f,r)}{r^k}.
\]
\end{definition}

\begin{theorem}[Sharp integer-order ray threshold]
\label{thm:generalray}
Let $k\ge1$ be an integer and let $A\subset(0,\infty)$ be discrete with
\begin{equation}
 N_A(r)=d r^k+O(r^{k-1})
 \label{eq:generalcount}
\end{equation}
for some $d>0$.  Then:
\begin{enumerate}
\item every nonzero entire $f$ vanishing on $A$ and having finite order-$k$ type satisfies
\[
 \tau_k(f)\ge\pi d;
\]
\item the product
\begin{equation}
 H_{A,k}(z)=z^{\epsilon}
 \prod_{a\in A}\left(1-\frac{z^{2k}}{a^{2k}}\right),
 \label{eq:generalproduct}
\end{equation}
where $\epsilon$ accounts for a possible zero at the origin, is entire and satisfies
\[
 \log M(H_{A,k},r)=\pi d\,r^k+O(r^{k-1}+\log r);
\]
\item therefore the least finite order-$k$ type of a nonzero entire function vanishing on
$A$ is exactly $\pi d$.
\end{enumerate}
\end{theorem}

\begin{proof}
For the lower bound, choose the principal $k$th root in the right half-plane and set
$F(w)=f(w^{1/k})$.  It is holomorphic there, has exponential type at most $\tau_k(f)$, and
has positive zeros $a^k$ with counting function $dT+O(T^{1-1/k})$.  Carleman's formula gives
$\tau_k(f)\ge\pi d$ exactly as in Theorem~\ref{thm:raydensity}.

For the upper bound, the product converges because $2k>k$.  Its exact maximum modulus is
\[
 r^{\epsilon}\prod_{a\in A}\left(1+\frac{r^{2k}}{a^{2k}}\right).
\]
Stieltjes integration gives the main term
\[
 2k d r^k\int_0^\infty\frac{u^{k-1}}{1+u^{2k}}\,du
 =2k d r^k\cdot\frac{\pi}{2k}=\pi d r^k.
\]
The counting error gives the stated lower-order term.
\end{proof}

\begin{corollary}[Higher-rank positive semigroups]
\label{cor:rankksemigroup}
Let $\alpha_1,\dots,\alpha_k>0$ be linearly independent over $\Q$ and define
\[
 \mathcal S_{\boldsymbol\alpha}
 =\left\{n_1\alpha_1+\cdots+n_k\alpha_k:n_i\in\Nzero\right\}.
\]
Then
\[
 N_{\boldsymbol\alpha}(R)
 =\frac{R^k}{k!\alpha_1\cdots\alpha_k}+O_{\boldsymbol\alpha}(R^{k-1}),
\]
and the least finite order-$k$ type of a nonzero entire function vanishing on this semigroup
is
\begin{equation}
 \frac{\pi}{k!\alpha_1\cdots\alpha_k}.
 \label{eq:rankkconstant}
\end{equation}
\end{corollary}

\begin{proof}
The counting formula is the standard lattice-point estimate for a simplex.  Apply
Theorem~\ref{thm:generalray}.
\end{proof}

For $k=2$, $(\alpha_1,\alpha_2)=(1,\beta)$, formula
\eqref{eq:rankkconstant} is exactly $\pi/(2\beta)$.  This general theorem also clarifies the
rotational order: at order $k$, the optimal genus-zero regularization uses $2k$ roots of
unity.

\begin{statusbox}{The general integer-order theorem and semigroup corollary are
\statusproved.}
\end{statusbox}

\section{Consequences for construction strategies}
\label{sec:strategy}

The exact threshold and translation-rank theorem reorganize the analytic frontier.  Several
classes of construction can now be rejected before any detailed coefficient calculation.

\subsection{Strategies now closed}

\paragraph{Vanishing correction terms.}
Every nonzero null function has type at least $\Cbeta$, over thirteen thousand times the
Pila threshold at the kernel lower bound $\beta=14$.  A ``build approximately, then repair
on all nodes'' strategy cannot be subcritical.

\paragraph{One fixed type-zero template.}
A Pila-small function can coincide with a fixed order-one or type-zero template on at most
$7.18\times10^{-5}$ of the full Carleman density when $\beta=14$.  Sparse deviations from
$2^z$, $3^z$, a polynomial, or an exponential polynomial are therefore impossible.

\paragraph{Small finite alphabets.}
Any fixed local translation stencil with fewer than about $13{,}933$ possible outputs at
$\beta=14$ either fails or becomes a global identity.  This rules out low-state carry
machines and bounded-defect decoders at Pila type.

\paragraph{Polynomial coordinate formulas.}
Any algebraic-coefficient polynomial in $(n,k)$ collapses to a constant below $\Cbeta$.
Thus a successful value table cannot be a low-degree arithmetic formula in the canonical
coordinates.

\paragraph{Finite linear recurrence rank.}
Two coordinate recurrences force type zero.  A successful positive-type table must be
non-C-finite in at least one direction and have unbounded Hankel complexity.

\subsection{Strategies that survive}

\paragraph{Direct nonholonomic interpolation.}
The value table must be constructed directly, without a low-type null correction.  It must
use an unbounded or very large alphabet and evade finite recurrence structure.  Candidates
include output-sensitive Newton schemes whose denominators and signs are coupled across both
coordinates, rather than fixed universal kernels.

\paragraph{Several algebraically independent interpolants.}
The multi-function gate permits larger types.  Instead of forcing one function below
$O(\beta^{-2})$, one may seek $q$ algebraically independent integer-valued functions near
the scale $O_q(\beta^{-1-1/q})$.  Their values need not encode the same observable.

\paragraph{One function with a rich translate orbit.}
A single function of normalized type below $A_0$ produces infinitely many integer-valued
translates.  Pila then forces finite transformal rank.  Proving that this finite-rank
alternative collapses to type zero would turn the existence of any such function into a
contradiction.

\paragraph{Nonlinear partial-difference geometry.}
The remaining relation should be attacked as a pair of commuting algebraic dynamics rather
than as another interpolation determinant.  Degree growth, Nevanlinna theory for differences,
and the classification of low-entropy birational maps are natural inputs.

\subsection{A concrete staged campaign}

\begin{enumerate}
\item \textbf{Formalize the exact zero threshold.}  This prevents the false infinite-type
claim from propagating and supplies a reusable uniqueness lemma.

\item \textbf{Build a two-shift difference-algebra package.}  Formalize translation fields,
relative transcendence degree, and extraction of a polynomial partial-difference equation
from Theorem~\ref{thm:translationrank}.

\item \textbf{Classify scalar first-order equations.}  Extend
Proposition~\ref{prop:polyiteration} from polynomials to rational and algebraic
correspondences, keeping explicit finite-order hypotheses.

\item \textbf{Prove a dynamical-degree lemma.}  Show that superlinear degree growth of a
shift correspondence forces infinite Nevanlinna order unless the entire curve lies in an
exceptional subvariety.

\item \textbf{Analyze the degree-one cases.}  Reduce commuting low-entropy maps to additive,
multiplicative, or elliptic group actions.  Apply the dense subgroup generated by $1$ and
$\beta$ and the integer-value condition.

\item \textbf{Parallel computational search.}  Search for finite value tables whose
interpolants have small empirical quadratic indicator and whose translate fields exhibit
large algebraic rank.  The purpose is not numerical evidence for the conjecture, but
identification of a viable nonholonomic architecture.
\end{enumerate}

\begin{target}[Translation-rank closure criterion]
\label{target:translationclosure}
Prove the following statement for some absolute $A>0$: if $g$ is entire,
$g(\Sbeta)\subseteq\Z$, $2\beta\tau_2(g)<A$, and the two-shift translation field has finite
relative transcendence degree over $K_0$, then $g$ is algebraic over $K_0$.

Any $A\ge A_0=29/75$ closes the multi-function route developed here.  Even a smaller
positive $A$ would improve the one-function landscape if combined with a construction at the
corresponding scale.
\end{target}

\section{Control tests and adversarial checks}
\label{sec:controls}

Every new mechanism was checked against the standard controls.

\subsection{Integer generators}

If $\beta$ is an integer, the representation $n+k\beta$ is not unique and the set
$\Nzero+\beta\Nzero$ is simply $\Nzero$.  Its counting function is linear, not quadratic.
Hence $d_\beta=1/(2\beta)$ is not applicable.  The exact threshold theorem relies on the
free rank-two monoid supplied by irrationality.  It does not falsely exclude the obvious
integer solutions.

\subsection{Arbitrary irrational generators}

The zero-set results themselves hold for every irrational $\beta>0$, independently of
$2^\beta$ and $3^\beta$.  What is arithmetic is the assertion that $2^s$ and $3^s$ are
integers for all $s\in\Sbeta$, which is used in the integer-valued spectrum and Pila
sections.  This separation is useful: the analytic theorems can be verified without assuming
any disputed arithmetic statement.

\subsection{No contradiction from the explicit null function}

The function $H_\beta$ is integer-valued on $\Sbeta$ because it is zero there, but its type
is $\Cbeta$, far outside Pila's small-growth range.  Pila's theorem therefore does not
produce an algebraic-dependence contradiction from $H_\beta$.  The construction corrects
the spectrum; it does not prove the conjecture by itself.

\subsection{No hidden cancellation in the type computation}

The maximum-modulus formula for $H_{m,\beta}$ is exact at every radius, not merely an upper
bound.  All factors attain equality on one common ray.  The type computation therefore
cannot be invalidated by phase cancellation between factors.

\subsection{The Carleman constant}

The factor $\pi$ is calibrated by the classical model of zeros at the positive integers:
a right-half-plane function of exponential type less than $\pi$ cannot vanish at every
nonnegative integer unless it is zero.  Under the squaring map, the monoid density becomes
$1/(2\beta)$, giving $\pi/(2\beta)$.  This provides an independent normalization check.

\subsection{Strict inequalities in Pila's theorem}

The exact numerical tables report suprema.  A proof must choose a slightly smaller type so
that the weighted-simplex inequality is strict.  This is a harmless open-margin issue, not a
claim of equality at the numerical boundary.

\begin{statusbox}{The control checks expose no contradiction or control-pair false positive.}
\end{statusbox}

\section{Literature audit and relation to current work}
\label{sec:literature}

The new results sit at the intersection of three established subjects: integer-valued entire
functions, uniqueness sets of prescribed angular density, and finite-order complex difference
equations.  This section records exactly which external facts are used and which recent
results merely provide context.

\subsection{Pila's integer-valued entire-function theorem}

The quantitative input is Pila's theorem on several entire functions sharing a set of
arguments of integrality~\cite{Pila2008}.  In the notation of that paper, the increasing
enumeration of
\[
 X_{\alpha,\beta}=\{i\alpha+j\beta:i,j\in\Nzero\}
\]
satisfies $x_n\sim(2\alpha\beta n)^{1/2}$.  Theorem~4.1 permits several additional
functions $g_i$ with normalized rates $\lambda_i$ and produces algebraic dependence when
\[
 \sum_{\mathbf j\cdot\boldsymbol\lambda\le A}
 (A-\mathbf j\cdot\boldsymbol\lambda)^k
 >s k!\,\omega_1\cdots\omega_k,
 \qquad
 A<\frac{\chi(X,t)}{t+(s-1)^{-1}}.
\]
Our specialization $k=t=s=2$, $(\alpha,\beta)=(1,\beta)$,
$(a,b)=(\log2,\log3)$ is literal.  The only addition in this report is the sharp promotion
from dependence on $\Sbeta$ to an identity of entire functions, using
Theorem~\ref{thm:sharpthreshold}, and the systematic application to a full translation
orbit.

The one-function constant $1/832$ is Pila's deliberately conservative published constant.
The exact finite sum $\Phi_1$ gives the slightly larger numerical values in
Section~\ref{sec:multipila}; this is not a strengthening of Schneider's method, only the
retention of the finite monomial sum before replacing it by an integral.

\subsection{Carlson, Carleman, and angular-density uniqueness}

Carleman's formula in a half-plane and Carlson's theorem are classical; a standard source
is Levin~\cite{Levin1980}.  A modern refinement of Carlson's theorem in terms of spectral
measures is due to Vagharshakyan~\cite{Vagharshakyan2021}.  Grishin and Sodin developed the
relation between growth on a ray, angular zero density, and uniqueness
sets~\cite{GrishinSodin1990}.

Recent work of Kononova gives a geometric theory of critical zero-set and uniqueness-set
types for sets with angular density~\cite{KononovaI,KononovaII,KononovaIII}.  In particular,
the order-two case is now treated explicitly from the viewpoint of type-minimizing angular
measures.  The present monoid provides a particularly elementary test case: after the map
$z\mapsto z^2$, its positive zeros have linear density $1/(2\beta)$, and fourfold rotational
completion gives the extremizer directly.  Thus our value $\pi/(2\beta)$ is consistent with
the general angular-density geometry, but the proof here requires neither a Lindelof
condition nor a completely-regular-growth classification.

\subsection{Difference Nevanlinna theory and partial difference algebra}

Difference analogues of the logarithmic derivative lemma and the use of finite order to
detect integrable difference equations were developed by Halburd and
Korhonen~\cite{HalburdKorhonen2006,HalburdKorhonen2007} and independently in related form by
Chiang and Feng~\cite{ChiangFeng2008}.  Those results strongly support the degree-growth
strategy in Section~\ref{sec:closedsubcases}: nonlinear recurrences of genuinely expanding
degree usually force excessive Nevanlinna growth.

They do not, however, provide the classification needed here.  Our relation can involve
several transforms, coefficients in $K_0=\C(2^z,3^z)$, two commuting irrationally related
shifts, and an algebraic variety of positive dimension.  Recent partial difference algebra
develops dimension polynomials and finite generation for several commuting operators, for
example McGrath~\cite{McGrath2025}; this supplies an appropriate language for the translation
field, but not yet an analytic finite-order classification.

\subsection{Adjacent transcendence developments}

Current work on matrices of logarithms, arithmetic holonomy, and additive relations in
irrational powers is relevant to the wider conjectural landscape
\cite{Dasgupta2023,CDT2025,Harrison2026}.  These works sharpen weight-one logarithmic
independence, effective Diophantine approximation, and additive expansion under irrational
powers.  None presently supplies the missing weight-two separation for
\[
 \log M\,\log3-\log A\,\log2=0,
\]
nor does any of them imply the two-shift classification target.  They are therefore recorded
as adjacent techniques rather than quoted as solutions.

\subsection{Evidence policy}

Only theorem statements in primary mathematical sources and proofs written in this report
are used.  Competitive announcements, press reports, and social-media claims are not
mathematical premises.  In particular, no claimed recent machine-generated breakthrough is
used unless a checkable proof or formal artifact is available.

\begin{statusbox}{The literature audit is contextual.  The external inputs actually used in
proofs are Pila's theorem and the classical Carleman formula; all other citations identify
nearby theories and possible tools.}
\end{statusbox}

\section{Formalization and verification roadmap}
\label{sec:formalization}

The new results divide cleanly into an elementary algebraic layer, a classical analytic
layer, and an external transcendence layer.  This separation is useful both for Lean and for
ordinary refereeing.

\subsection{Layer A: finite sums, products, and zero covers}

The following statements require no transcendence theory.

\begin{enumerate}
\item The exact floor sum in Proposition~\ref{prop:counting} and the estimate
$N_\beta(R)=R^2/(2\beta)+O(R)$.
\item Convergence of $H_{m,\beta}$ from $\sum s^{-m}<\infty$.
\item The exact maximum-modulus identity in Lemma~\ref{lem:maxmod}.
\item The beta-integral evaluation
\[
 \int_0^\infty\frac{u}{1+u^m}\,du
 =\frac\pi m\csc\frac{2\pi}{m}.
\]
\item Finite zero-cover, template, and alphabet deductions once the Carleman-density lemma is
available.
\item The finite-difference polynomial decoder and the constant-increment arguments.
\item Exact evaluation of $\Phi_q$ and the tables in Sections~\ref{sec:multipila} and
\ref{sec:translationrank}.
\end{enumerate}

The main formalization choices are representations of the locally finite monoid and a
library interface for infinite products indexed by a countable discrete subset of
$\R_{>0}$.  One can initially avoid asymptotic Stieltjes integration by proving two-sided
Riemann-sum bounds with explicit $O(R)$ counting errors.

\subsection{Layer B: the Carleman bridge}

The analytic kernel is the following reusable theorem.

\begin{target}[Formal Carleman-density kernel]
For a function holomorphic in the right half-plane with
$\log^+|F(w)|\le(\sigma+o(1))|w|$, formalize
\[
 \limsup_{R\to\infty}\frac1{\log R}
 \sum_{\substack{F(a)=0\\1<a<R}}
 \left(\frac1a-\frac a{R^2}\right)
 \le\frac\sigma\pi
\]
for positive real zeros counted with multiplicity.
\end{target}

Once this theorem exists, the sharp lower bound follows from the principal square-root map
in a few lines.  A pragmatic first formalization may state Carleman's formula as a library
theorem and separately verify that its boundary terms have the needed bounds.  A fully
foundational proof can follow Appendix~\ref{app:carleman} through Green's identity on a
half-annulus with small disks removed around zeros.

\subsection{Layer C: Pila as an explicit external theorem}

Pila's theorem is far beyond the reasonable scope of a first formalization.  It should be
encapsulated with all hypotheses visible:

\begin{itemize}
\item the scale property and the constant $\chi(X,2)$;
\item the normalized growth rates $\omega_i$ and $\lambda_j$;
\item integrality on the ordered scale;
\item the weighted-simplex strict inequality;
\item the conclusion of polynomial dependence on the scale, together with the weighted
degree support needed for the growth promotion.
\end{itemize}

The promotion from a relation on $\Sbeta$ to an entire identity belongs to Layer B and can
be kernel-checked independently.  This prevents a hidden strengthening of the cited theorem.

\subsection{Suggested theorem dependency order}

\begin{enumerate}
\item \texttt{ConditionalMonoid.counting}.
\item \texttt{ConditionalMonoid.reciprocalPowerSummable}.
\item \texttt{RotationalProduct.converges}.
\item \texttt{RotationalProduct.maxModulus}.
\item \texttt{RotationalProduct.typeFormula}.
\item \texttt{Carleman.positiveRayDensityBound}.
\item \texttt{ConditionalMonoid.sharpVanishingType}.
\item \texttt{ConditionalMonoid.finiteZeroCover}.
\item \texttt{ConditionalMonoid.polynomialDecoderRigidity}.
\item \texttt{PilaMulti.specializedGate} as an external-input wrapper.
\item \texttt{TranslationOrbit.finiteRelativeTranscendenceDegree}.
\end{enumerate}

\subsection{Falsification tests}

Every formal or paper implementation should include the following controls.

\begin{longtable}{@{}p{0.31\textwidth}p{0.29\textwidth}p{0.32\textwidth}@{}}
\toprule
Test & Expected result & Failure it detects\\
\midrule
\endfirsthead
\toprule
Test & Expected result & Failure it detects\\
\midrule
\endhead
$A=\N$ in Theorem~\ref{thm:generalray} & minimum exponential type $\pi$ & missing factor $\pi$ or $2$ in Carleman\\
$H_{4,\beta}$ on $z=r e^{i\pi/4}$ & exact product maximum & unjustified upper-bound-to-type equality\\
$m=3,4,5$ in \eqref{eq:Hmtype} & unique minimum at $m=4$ & wrong beta-integral exponent\\
Integer $\beta$ & reject free-monoid counting & accidental application to nonunique coordinates\\
Zero function in a zero cover & remove before applying type bound & division by a nonexistent nonzero function\\
$\lambda\downarrow0$ in $\Phi_q$ & divergence like $\lambda^{-q}$ & missing high-degree monomials\\
Control pair $\beta=n\in\N$ in Pila specialization & hypotheses of counterexample package fail & false contradiction with obvious integer exponents\\
One periodic shift only & positive-type examples remain possible & overclaiming from one difference equation\\
\bottomrule
\end{longtable}

\begin{statusbox}{The dependency map and controls are proposed formalization architecture,
not claims that a new Lean build has been completed.}
\end{statusbox}

\section{Status matrix for this report}
\label{sec:status}

\begin{longtable}{@{}p{0.43\textwidth}p{0.18\textwidth}p{0.31\textwidth}@{}}
\toprule
Statement & Status & Dependency\\
\midrule
\endfirsthead
\toprule
Statement & Status & Dependency\\
\midrule
\endhead
Conditional monoid count $R^2/(2\beta)+O(R)$ & \statusproved & elementary floor sum\\
Exact-zero genus-two product has infinite quadratic type & \statusproved & partial summation\\
$H_{m,\beta}$ converges and has exact maximum modulus & \statusproved & infinite-product criterion\\
$\tau_2(H_{m,\beta})=\pi/(2\beta\sin(2\pi/m))$ & \statusproved & Stieltjes integration, beta integral\\
Predecessor's infinite-type proposition & \statuscorrection & contradicted by $H_{4,\beta}$\\
Sharp lower bound $\tau_2\ge\pi/(2\beta)$ & \statusproved & classical Carleman formula\\
Complete vanishing spectrum $[\pi/(2\beta),\infty)$ & \statusproved & sharp bound and rotated Gaussian\\
Integer-valued spectrum contains $\{0\}\cup[\pi/(2\beta),\infty)$ & \statusproved & $2^z+H_{\beta,a}$\\
Pila exclusion through $1/(832\beta^2\log2\log3)$ & \statusclassical & Pila plus inherited growth lemma\\
Finite zero-cover and alphabet inequalities & \statusproved & Carleman density\\
Polynomial decoder rigidity & \statusproved & subcritical uniqueness, two differences\\
Two constant recurrences imply type zero & \statusproved & Fourier--Floquet decomposition\\
Weighted-simplex multi-function gate & \statusproved & Pila Theorem~4.1 plus sharp promotion\\
Finite translation transcendence degree & \statusproved & multi-function gate\\
Two-shift finite-order classification & \statusopen & new principal target\\
Dynamical-degree reduction & \statusopen & proposed intermediate theorem\\
Full Alaoglu--Erd\H{o}s conjecture & \statusopen & not proved here\\
\bottomrule
\end{longtable}

\section{Conclusions and the exact new frontier}
\label{sec:conclusion}

The investigation produced a genuine correction and a new route, but not a proof of the
conjecture.  The correction is mathematically decisive: the zero-set branch does not fail by
infinite quadratic type.  It has an exact, finite critical constant
\[
 \boxed{\Cbeta=\frac{\pi}{2\beta}}.
\]
This number is simultaneously

\begin{itemize}
\item the least quadratic type of a nonzero entire function vanishing on the conditional
solution monoid;
\item the threshold below which interpolation on that monoid is unique;
\item the total type budget required by any finite cover of the monoid by zero sets;
\item the state-complexity constant governing finite alphabets and templates;
\item the order-two instance of the general ray-density constant $\pi d$.
\end{itemize}

The explicit extremizer
\[
 H_\beta(z)=z\prod_{s\in\Sbeta\setminus\{0\}}
                 \left(1-\frac{z^4}{s^4}\right)
\]
shows exactly how the earlier argument failed: at integer order, adding three rotated copies
of the positive zero ray cancels the critical angular moment.  The accompanying Carleman
argument proves that this cancellation is optimal.

The exact threshold does not make Pila's one-function construction easy.  The quantitative
gap remains
\[
 \frac{\Cbeta}{\TP}
 =416\pi\beta\log2\log3,
\]
which is already about $1.39\times10^4$ at $\beta=14$.  It does, however, turn vague
statements about ``too many zeros'' into sharp no-go theorems.  Low-type null corrections,
small finite alphabets, polynomial coordinate decoders, and finite recurrence rank are now
closed with exact constants.

The most promising new consequence is the translation-orbit theorem.  Pila's
multi-function result converts any positive-type integer-valued interpolant with
$2\beta\tau_2(g)<29/75$ into a function whose entire two-dimensional translation orbit has
finite relative transcendence degree over $\C(2^z,3^z)$.  Therefore the next analytic
question is not to optimize yet another universal determinant.  It is to classify
finite-order entire solutions of algebraic partial-difference equations under two commuting
real shifts whose ratio is irrational.

A closure theorem of the following form would be sufficient:
\[
 \begin{gathered}
 g(\Sbeta)\subseteq\Z,\qquad
 0<2\beta\tau_2(g)<29/75,\\
 \trdeg_{\C(2^z,3^z)}
 \C(2^z,3^z)(g(z+n+k\beta):n,k\ge0)<\infty
 \end{gathered}
 \quad\Longrightarrow\quad
 g\text{ is algebraic over }\C(2^z,3^z).
\]
The linear, periodic, polynomial-decoder, and scalar expanding-polynomial cases are already
closed in this report.  The surviving case is nonlinear but highly structured: finite
transformal rank, two dense real directions, finite Nevanlinna order, and integer values on
the positive orbit.  That package is substantially narrower than the original interpolation
problem and is the recommended focus for the next report.

\begin{statusbox}{The full conjecture remains open.  The durable progress is the corrected
sharp type theorem, its complexity consequences, the multi-function gate, and the finite
translation-rank reduction.}
\end{statusbox}

\appendix

\section{Derivation of the half-plane Carleman inequality}
\label{app:carleman}

For completeness we derive the constant used in Lemma~\ref{lem:carleman}.  The purpose of
the calculation is not to reproduce every regularization detail, but to make the coefficient
$1/\pi$ auditable.

Let
\[
 D_R^+=\{z:|z|<R,\ \Re z>0\}
\]
and suppose first that $F$ is holomorphic in a neighborhood of its closure, has no zeros on
the boundary, and $F(0)\ne0$.  Apply Green's identity to $\log|F(z)|$ and the harmonic
weight
\[
 v_R(z)=\Re\left(\frac1z-\frac z{R^2}\right).
\]
The weight vanishes on the semicircle $|z|=R$, is positive inside the right half-disk, and
has normal derivative giving the standard Carleman kernel.  Removing small circles around
the zeros $a_j$ and around the origin, then letting their radii tend to zero, yields
\begin{align}
 \sum_{|a_j|<R}
 \left(\Re\frac1{a_j}-\frac{\Re a_j}{R^2}\right)
 &=\frac1{2\pi}\int_{-R}^{R}
   \log|F(iy)|\left(\frac1{y^2}-\frac1{R^2}\right)\,dy
 \label{eq:fullcarleman}\\
 &\quad+\frac1{\pi R}\int_{-\pi/2}^{\pi/2}
   \log|F(Re^{i\theta})|\cos\theta\,d\theta
 +O_F(1).
 \nonumber
\end{align}
Zeros are counted with multiplicity.  The bounded term contains the contribution near the
origin and a normalization by $F(0)$; neither affects a coefficient of $\log R$.

For a positive real zero $a_j$, its summand is exactly
\[
 \frac1{a_j}-\frac{a_j}{R^2}.
\]
Discarding all other zero summands is legitimate because the full left side is nonnegative
term by term for zeros in $D_R^+$:
\[
 \Re\left(\frac1a-\frac a{R^2}\right)
 =\Re a\left(\frac1{|a|^2}-\frac1{R^2}\right)>0.
\]

Assume now that for every $\varepsilon>0$,
\[
 \log^+|F(z)|\le(\sigma+\varepsilon)|z|+O_\varepsilon(1)
 \qquad(\Re z\ge0).
\]
On the semicircle, the second integral in \eqref{eq:fullcarleman} is $O(1)$ after division by
$R$.  On the imaginary axis, symmetry and the positive part give
\begin{align*}
 &\frac1{2\pi}\int_{-R}^{R}
   \log^+|F(iy)|\left(\frac1{y^2}-\frac1{R^2}\right)\,dy\\
 &\qquad\le
 \frac{\sigma+\varepsilon}{\pi}
 \int_1^R\left(\frac1y-\frac y{R^2}\right)\,dy+O(1)\\
 &\qquad=
 \frac{\sigma+\varepsilon}{\pi}\log R+O(1).
\end{align*}
This proves \eqref{eq:carlemanineq}.  Boundary zeros and functions defined only in the open
half-plane are handled by applying the formula to $F(z+\delta)$ at regular radii and then
letting $\delta\downarrow0$.  Zeros at the origin are removed by dividing by a power of
$z$.

Two calibration checks are immediate.

\begin{enumerate}
\item For positive integer zeros, the left side is $\log R+O(1)$, so $\sigma\ge\pi$.
This is Carlson's constant.
\item For positive zeros of linear density $d$, the left side is $d\log R+O(1)$, so
$\sigma\ge\pi d$.
\end{enumerate}

Under $F(w)=f(\sqrt w)$, a quadratic-type bound $\tau_2(f)$ becomes the exponential-type
bound $\sigma=\tau_2(f)$, while the squared monoid has density $d=1/(2\beta)$.  The constant
$\pi/(2\beta)$ follows with no additional factor.

\section{Detailed asymptotics for the rotational product}
\label{app:productasymptotic}

This appendix makes the error in Theorem~\ref{thm:Hmtype} explicit.  Write
\[
 N_\beta(t)=d_\beta t^2+E_\beta(t),
 \qquad |E_\beta(t)|\le C_\beta'(t+1)
\]
for $t\ge1$.  From integration by parts,
\[
 \log M(H_{m,\beta},r)-\log r
 =m r^m\int_{s_0}^{\infty}
       \frac{N_\beta(t)}{t(t^m+r^m)}\,dt,
\]
where $s_0=\min(\Sbeta\setminus\{0\})$.

The main term is
\[
 m d_\beta r^2\int_0^\infty\frac{u}{1+u^m}\,du
 =\pi d_\beta\csc(2\pi/m)r^2.
\]
For the $O(t)$ portion of the error, split at $t=r$:
\begin{align*}
 r^m\int_{s_0}^{r}\frac{dt}{t^m+r^m}&=O(r),\\
 r^m\int_r^\infty\frac{dt}{t^m+r^m}&=O(r).
\end{align*}
For the bounded portion,
\begin{align*}
 r^m\int_{s_0}^{r}\frac{dt}{t(t^m+r^m)}&=O(\log r),\\
 r^m\int_r^\infty\frac{dt}{t(t^m+r^m)}&=O(1).
\end{align*}
This proves the stated $O_{\beta,m}(r+\log r)$ error.  The exact maximum-modulus ray turns
this asymptotic into an equality for the type, rather than only an upper estimate.

The beta integral follows by $v=u^m$:
\[
 \int_0^\infty\frac{u}{1+u^m}\,du
 =\frac1m\int_0^\infty\frac{v^{2/m-1}}{1+v}\,dv
 =\frac\pi m\csc\frac{2\pi}{m},
\]
using Euler's reflection formula.  At $m=4$ this is $\pi/4$, and the prefactor $m d_\beta$
gives $\pi d_\beta$.

\section{Details on two incommensurable difference equations}
\label{app:difference}

We record a slightly more invariant proof of Lemma~\ref{lem:twoCfinite}.  Let $T_h$ denote
translation by $h$.  If $A(T_1)G=0$, decompose the finite-dimensional primary spectrum of
$T_1$ according to the roots $\mu$ of $A$.  Choosing one logarithm $\lambda$ with
$e^\lambda=\mu$, the generalized $\mu$-component has the form
\[
 e^{\lambda z}\sum_{j=0}^{m_\mu-1}z^jP_j(z),
\]
where $P_j$ is entire and $1$-periodic.  Every entire $1$-periodic function is the pullback
of a holomorphic function on $\C^*$ and has a locally normally convergent Laurent expansion
\[
 P_j(z)=\sum_{n\in\Z}c_{j,n}e^{2\pi i n z}.
\]

On the finite polynomial jet attached to the Fourier mode $n$, $T_\beta$ has generalized
eigenvalue
\[
 \nu_{\lambda,n}=e^{\beta(\lambda+2\pi i n)}.
\]
For fixed $\lambda$, these values are pairwise distinct because $\beta\notin\Q$.  Between
two different choices $\lambda_1,\lambda_2$, a collision
$\nu_{\lambda_1,n}=\nu_{\lambda_2,m}$ can occur for at most one difference $n-m$; two such
collisions would imply an integer multiple of $\beta$ is an integer.  Therefore every
$T_\beta$ generalized eigenspace receives only finitely many Fourier modes from the finite
set of $T_1$ branches.

The equation $B(T_\beta)G=0$ permits only generalized eigenvalues among the finitely many
roots of $B$, with bounded jet multiplicity.  Hence only finitely many Fourier modes remain,
and $G$ is an exponential polynomial.  The same proof shows a useful extension: if the two
right sides lie in finite-dimensional translation-invariant spaces, then $G$ lies in such a
space as well.

\section{Reproducibility code for constants and tables}
\label{app:code}

The following pure-Python script reproduces the constants, equal-type thresholds, and
translation-rank table.  It uses only integer binomial coefficients and double-precision
transcendental constants; replacing the final bisection by rational interval arithmetic is
straightforward.

\begin{lstlisting}[style=pythonstyle,caption={Reproduction of the numerical constants}]
from __future__ import annotations

import math
from math import comb

A0 = 29.0 / 75.0
LOG_PRODUCT = math.log(2.0) * math.log(3.0)


def phi(q: int, A: float, lam: float) -> float:
    # Exact finite sum Phi_q(A, lambda) for positive lambda.
    if q <= 0:
        raise ValueError('q must be positive')
    if lam <= 0.0:
        return math.inf
    degree = math.floor(A / lam + 1.0e-14)
    return sum(
        comb(ell + q - 1, q - 1) * (A - ell * lam) ** 2
        for ell in range(degree + 1)
    )


def common_type_threshold(q: int, beta: float) -> float:
    # Supremal tau for which Phi_q exceeds Pila's right side.
    rhs = 8.0 * beta * LOG_PRODUCT
    lo = 0.0
    hi = A0 / (2.0 * beta)
    for _ in range(160):
        mid = (lo + hi) / 2.0
        if phi(q, A0, 2.0 * beta * mid) > rhs:
            lo = mid
        else:
            hi = mid
    return lo


def q_star(beta: float, lambda_ratio: float) -> int:
    # Least q with Phi_q(A0, lambda_ratio*A0) above Pila's right side.
    if not (0.0 < lambda_ratio < 1.0):
        raise ValueError('lambda_ratio must lie in (0,1)')
    rhs = 8.0 * beta * LOG_PRODUCT
    lam = lambda_ratio * A0
    q = 1
    while phi(q, A0, lam) <= rhs:
        q += 1
    return q


def report() -> None:
    chi = (1.0 - math.sqrt(2.0) + math.log(2.0)
           - math.log(math.sqrt(2.0) - 1.0))
    gap_slope = 416.0 * math.pi * LOG_PRODUCT
    print(f'chi_2(2) = {chi:.12f}')
    print(f'A0 = {A0:.12f}; A0/chi_2(2) = {A0/chi:.12f}')
    print(f'C_beta/T_P per beta = {gap_slope:.12f}')

    for beta in (14.0, 27.0):
        print(f'beta = {beta:g}')
        for q in (1, 2, 3, 4, 5, 8):
            print(q, f'{common_type_threshold(q, beta):.12e}')
        values = [q_star(beta, r) for r in (0.05, 0.10, 0.25, 0.50)]
        print('q*:', values)


if __name__ == '__main__':
    report()
\end{lstlisting}

The expected leading output is
\begin{verbatim}
chi_2(2) = 1.160307205206
A0 = 0.386666666667; A0/chi_2(2) = 0.333245079348
C_beta/T_P per beta = 995.206300791...
\end{verbatim}
The table values then agree with Sections~\ref{sec:multipila} and
\ref{sec:translationrank} to the printed digits.

\section{Dependency graph and next-proof checklist}
\label{app:dependency}

The logical flow of the report is
\[
 \boxed{\text{conditional monoid}}
 \Longrightarrow
 \boxed{N_\beta(R)=R^2/(2\beta)+O(R)}
 \Longrightarrow
 \begin{cases}
  \boxed{H_{4,\beta}\text{ gives type }\pi/(2\beta)},\\
  \boxed{\text{Carleman gives type }\ge\pi/(2\beta)}
 \end{cases}
\]
\[
 \Longrightarrow
 \boxed{\text{exact subcritical uniqueness}}
 \Longrightarrow
 \begin{cases}
  \boxed{\text{zero-cover/alphabet/decoder rigidity}},\\
  \boxed{\text{promotion of Pila relations to entire identities}}
 \end{cases}
\]
\[
 \Longrightarrow
 \boxed{\text{finite algebraic rank of the translate orbit}}
 \Longrightarrow
 \boxed{\text{two-shift difference-algebra target}}.
\]

A next proof attempt should answer the following questions in order.

\begin{enumerate}
\item Can every finite-order algebraic relation among transforms be replaced by one that is
minimal and irreducible under both shifts?
\item Does minimality force the induced correspondences on the transform variety to be
dominant and generically finite?
\item Can Nevanlinna growth exclude first dynamical degree greater than one?
\item Do commuting degree-one correspondences reduce, after a finite cover, to algebraic
group translations?
\item Can integer values on the dense positive orbit eliminate additive, multiplicative,
and elliptic nonconstant branches?
\item If a branch remains, does it construct an additional exponential $e^{cz}$ with
algebraic values at both $1$ and $\beta$, thereby invoking six exponentials?
\end{enumerate}

The final question is especially attractive: a classification need not directly prove that
$g$ is algebraic over $K_0$.  It is enough to extract one genuinely new order-one
exponential direction whose values at both generators are algebraic.  That would convert the
two-shift geometry back into a three-base six-exponentials contradiction.

\begin{thebibliography}{99}

\bibitem{Unified6}
V.~Reshetnikov and the ProveIt project,
\emph{The Alaoglu--Erd\H{o}s Integer-Exponent Conjecture: Machine-Checked Partial Results,
Exact Conditional Structure, Determinant Methods, and the Remaining Barrier},
unified report, August 22, 2026; source supplied as
\texttt{alaoglu-erdos-unified-report-6.zip}.

\bibitem{Pila2008}
J.~Pila,
\emph{Entire functions sharing arguments of integrality, II},
Publ. Math. Debrecen \textbf{73} (2008), no.~1--2, 215--231,
\href{https://doi.org/10.5486/PMD.2008.4212}{doi:10.5486/PMD.2008.4212}.

\bibitem{Levin1980}
B.~Ya. Levin,
\emph{Distribution of Zeros of Entire Functions}, revised edition,
Translations of Mathematical Monographs, vol.~5,
American Mathematical Society, Providence, 1980.

\bibitem{Vagharshakyan2021}
A.~Vagharshakyan,
\emph{A refinement of Carlson's theorem},
\href{https://arxiv.org/abs/2108.12846}{arXiv:2108.12846}, 2021.

\bibitem{GrishinSodin1990}
A.~F. Grishin and M.~L. Sodin,
\emph{Growth along a ray, distribution of the zeros of an entire function of finite order
with respect to the argument, and a uniqueness theorem},
J. Soviet Math. \textbf{49} (1990), no.~6, 1269--1279.

\bibitem{KononovaI}
A.~Kononova,
\emph{Uniqueness sets with angular density for spaces of entire functions, I: Basics},
\href{https://arxiv.org/abs/2503.07225}{arXiv:2503.07225}, 2025; revised 2026.

\bibitem{KononovaII}
A.~Kononova,
\emph{Uniqueness sets with angular density for spaces of entire functions, II: $\rho=1$},
\href{https://arxiv.org/abs/2504.08116}{arXiv:2504.08116}, 2025.

\bibitem{KononovaIII}
A.~Kononova,
\emph{Uniqueness sets with angular density for spaces of entire functions, III: how to
minimize the type},
\href{https://arxiv.org/abs/2605.20380}{arXiv:2605.20380}, 2026.

\bibitem{HalburdKorhonen2006}
R.~G. Halburd and R.~J. Korhonen,
\emph{Difference analogue of the lemma on the logarithmic derivative with applications to
difference equations},
J. Math. Anal. Appl. \textbf{314} (2006), no.~2, 477--487,
\href{https://doi.org/10.1016/j.jmaa.2005.04.010}{doi:10.1016/j.jmaa.2005.04.010}.

\bibitem{HalburdKorhonen2007}
R.~G. Halburd and R.~J. Korhonen,
\emph{Finite-order meromorphic solutions and the discrete Painleve equations},
Proc. Lond. Math. Soc. (3) \textbf{94} (2007), no.~2, 443--474,
\href{https://doi.org/10.1112/plms/pdl012}{doi:10.1112/plms/pdl012}.

\bibitem{ChiangFeng2008}
Y.-M. Chiang and S.-J. Feng,
\emph{On the Nevanlinna characteristic of $f(z+\eta)$ and difference equations in the
complex plane},
Ramanujan J. \textbf{16} (2008), 105--129,
\href{https://doi.org/10.1007/s11139-007-9101-1}{doi:10.1007/s11139-007-9101-1}.

\bibitem{McGrath2025}
O.~McGrath,
\emph{Dimension polynomials for affine partial difference algebraic groups},
\href{https://arxiv.org/abs/2511.20350}{arXiv:2511.20350}, 2025.

\bibitem{Dasgupta2023}
S.~Dasgupta,
\emph{Ranks of matrices of logarithms of algebraic numbers I: the theorems of Baker and
Waldschmidt--Masser},
\href{https://arxiv.org/abs/2303.02037}{arXiv:2303.02037}, 2023.

\bibitem{CDT2025}
F.~Calegari, V.~Dimitrov, and Y.~Tang,
\emph{Arithmetic holonomy bounds and effective Diophantine approximation},
\href{https://arxiv.org/abs/2510.04156}{arXiv:2510.04156}, 2025.

\bibitem{Harrison2026}
J.~Harrison,
\emph{Additive relations in irrational powers},
\href{https://arxiv.org/abs/2512.04081}{arXiv:2512.04081}, version~3, 2026.

\end{thebibliography}

\end{document}
